\documentclass[11pt,a4paper]{article}
\usepackage{amsmath,amssymb,amsthm,mathtools}
\usepackage{mathpazo}
\usepackage[T1]{fontenc}
\usepackage[utf8]{inputenc}
\renewcommand{\ttdefault}{pcr}
\usepackage[margin=2.7cm]{geometry}
\usepackage{booktabs,array,graphicx}
\usepackage{enumitem}
\usepackage[round]{natbib}
\usepackage[colorlinks=true,linkcolor=blue!50!black,citecolor=blue!50!black,urlcolor=blue!50!black]{hyperref}
\usepackage{xcolor}
\setlength{\parskip}{0.35em}
\setlist{itemsep=0.2em,topsep=0.3em}
\setlength{\emergencystretch}{3em}

\newcommand{\E}{\mathbb{E}}
\newcommand{\PP}{\mathbb{P}}
\newcommand{\ind}{\mathbf{1}}
\newcommand{\X}{\mathcal{X}}
\newcommand{\A}{\mathcal{A}}
\newcommand{\dd}{\,\mathrm{d}}
\newcommand{\Pac}{P_{\mathrm{ac}}}
\newcommand{\Ps}{P_{\mathrm{s}}}
\newcommand{\thac}{\theta_{\mathrm{ac}}}
\newcommand{\ths}{\theta_{\mathrm{s}}}
\newcommand{\Cov}{\operatorname{Cov}}
\newcommand{\Var}{\operatorname{Var}}
\newcommand{\odds}{\operatorname{odds}}
\newcommand{\neff}{n_{\mathrm{eff}}}
\newcommand{\vacio}{\varnothing}
\usepackage[english]{babel}

\theoremstyle{plain}
\newtheorem{theorem}{Theorem}
\newtheorem{proposition}{Proposition}
\newtheorem{lemma}{Lemma}
\newtheorem{corollary}{Corollary}
\theoremstyle{definition}
\newtheorem{example}{Example}
\newtheorem{remark}{Remark}

\title{When channels measure differently: counterfactual evaluation of channel switching via change of measure}
\author{Pablo}
\date{October 24, 2026}

\begin{document}
\maketitle

\begin{abstract}
We estimate what would happen if a group of units (customers, patients, calls) were served by a channel other than the one serving them today, without a randomized experiment. We cast identification as a change of measure: the counterfactual mean is an integral over the donor population weighted by the Radon--Nikodym derivative, which a calibrated classifier identifies up to a constant. The focus is a problem that is common in operational data and rarely formalized: each channel records covariates in its own way. A variable recorded in only one channel breaks support by construction, and a confident classifier ($p=1$) reveals it; dropping it is safe when it does not predict the outcome given the other covariates. A variable measured with different instruments in each channel is more dangerous: it can preserve support, and the estimator then converges to a shifted value that no analysis of the observed data can distinguish from the truth. We also show that, without support, the doubly robust estimator implicitly extrapolates the outcome model, and we adapt sensitivity bounds for unobserved confounding. When the target group is a cutoff on a known score, the weights follow from the ordinary channel propensity with the cutoff imposed; estimating it with all units of the premium channel reduces the dispersion by up to 59\,\% in our simulation when overlap is good. A simulation with predictions registered before running it tests the main results and also finds, in our design, that the unit-level bootstrap overstates dispersion when the estimated weights are noisy. With public data (self-reported versus measured body mass index in NHANES), the predicted bias of differential measurement is approximately confirmed, while the overlap diagnostics improve rather than flag it.
\end{abstract}

\section{Introduction}\label{sec:intro}

Organizations allocate their users across channels of different quality and cost: a specialized service center versus a general one, a hospital versus primary care, a personal advisor versus an app. Before reallocating, they want to know how much is lost by moving a group to the basic channel and how much is gained by moving another group to the premium one. A randomized experiment would answer this, but it is often expensive, slow or infeasible; what is available are observational data from each channel.

The central question is always the same: \emph{what would the mean outcome of the target group be if the other channel served it?} That outcome is not observed. What is observed is the outcome of the units the other channel serves today, which are different. Reweighting them to resemble the target group is, formally, a change of measure.

The estimand coincides with the counterfactual mean of an effect on the treated, and the odds weights $p/(1-p)$ are well known \citep{hirano2003}. Estimating density ratios with classifiers is also well known \citep{qin1998,bickel2009,sugiyama2012}. We do not claim those ideas. Our contributions are:
\begin{enumerate}
  \item \textbf{Channel variables and differential measurement.} We treat covariates recorded in only one channel as the limiting case of the missingness-pattern approach \citep{dagostino2000,blake2020}, in which the pattern is determined by the treatment itself. Positivity then fails in the observed pattern, with an unsupported mass at least equal to the fraction of the target group with data, and the condition of \citet{blake2020} for the missing pattern holds, in particular, when the variable does not predict the outcome given the others (Propositions~\ref{prop:breaks} and~\ref{prop:drop}). We generalize this to variables measured with different instruments in each channel (Theorem~\ref{thm:measurement}): support can be preserved, and the estimator converges to a shifted value that the observed data cannot distinguish from the truth.
  \item \textbf{Support failure read through the classifier.} The unsupported part of $P$ is $P$ restricted to $\{p=1\}$. This is von Neumann's classical construction of the Lebesgue decomposition \citep[Theorem 6.10]{rudin1987}, which the overlap literature uses informally \citep{damour2021}; what we add is its reading in terms of the classifier and its consequences: the weighting estimator keeps the supported part, whereas the doubly robust estimator implicitly extrapolates the outcome model (Theorem~\ref{thm:lebesgue} and Corollary~\ref{cor:drnosupport}).
  \item \textbf{Known-cutoff weights.} When the target group is a cutoff on a known score, the weights are the ordinary channel propensity with the cutoff imposed (Proposition~\ref{prop:cutoff}), as in the usual practice of subgroup analysis \citep{rassen2012}. What is new is showing what is lost with the ``target versus donors'' classifier: estimating the propensity with all units of the premium channel greatly reduces the sign reversal of trimming produced by the classifier and, with good overlap, the dispersion, by up to 59\,\% in small samples; with medium or low overlap it does not reduce it (Section~\ref{sec:res-cutoff}).
  \item \textbf{Simulation with registered predictions.} The limit of each estimator is computed before running the simulation, and predictions added after seeing a result were registered before the scenario that tests them. We also report those that failed; among them, that the unit-level bootstrap, re-estimating the auxiliary models, overstates the dispersion by 35 to 60\,\% with noisy weights and is nearly calibrated with the known-cutoff weights (Section~\ref{sec:res-bootstrap}). With public data, the bias predicted for differential measurement is approximately confirmed, and the overlap diagnostics improve rather than flag it (Section~\ref{sec:illustrations}).
\end{enumerate}

As auxiliary results, we give the exact bias of capping the true weights with normalization (Proposition~\ref{prop:trimming}), which complements \citet{ma2020}, and we adapt the bounds of the marginal sensitivity model \citep{tan2006,dorn2023} to this estimand, with an exact multiplicative form for rare outcomes and a calibration through the channel model (Section~\ref{sec:sensitivity}). Trimming with estimated weights is developed in a companion note.

We use the language of measure theory (Radon--Nikodym derivatives) rather than density ratios. This is not a contribution but the standard formulation; it handles mixed covariates (continuous, categorical and with missing values) without assuming that densities exist.

\paragraph{Related work.} The estimand is the counterfactual mean of an effect on the treated, with the odds weights of \citet{hirano2003} and the doubly robust theory with cross-fitting of \citet{chernozhukov2018}. The same estimand, odds weights and doubly robust influence function appear in the literature on generalizing and transporting trial results to a target population \citep{westreich2017,dahabreh2019}. For limited overlap, \citet{crump2009} redefine the estimand on the region with good overlap. \citet{ma2020} characterize the bias of trimming by discarding in the weighting estimator, propose a correction and show that with heavy-tailed weights the limit may be non-normal; \citet{ma2023} extend the correction to doubly robust estimators under weak overlap. \citet{liu2023} compare trimming, truncation and overlap weights for the effect on the treated. \citet{manamkoong2025} split the estimand into a part with overlap and a part without it, and extrapolate the latter under a smoothness assumption.

\citet{heckman1997} found empirically that matching removes much of the bias when the comparison group comes from the same market and answers the same questionnaire as participants; Theorem~\ref{thm:measurement} formalizes that lesson as an indistinguishability result. \citet{cook2008} and \citet{steiner2011} also document the importance of measuring covariates in the same way and reliably, and that differential measurement error cannot be corrected without external information is known in the measurement-error literature; \citet{imai2010} study nonparametric identification and sensitivity analysis under differential measurement error of the treatment. On partially observed covariates, the missingness-pattern approach \citep{dagostino2000,blake2020} requires the variable not to confound where it is missing; our channel variables are the case in which the treatment fixes the pattern, and the condition we use is an outcome-based covariate selection criterion \citep{deluna2011,vanderweele2011}. In subgroup analysis, \citet{rassen2012} and \citet{dong2020} compare the propensity estimated in the full cohort with the one estimated within each subgroup. On diagnostics, \citet{westreich2011} show that a high AUC of the propensity model is neither necessary nor sufficient to control confounding; our simulation adds that it does not measure overlap either. For diagnosing positivity violations, see \citet{petersen2012} and \citet{damour2021}.
\section{Setup}\label{sec:setup}

We work on a probability space $(\Omega,\mathcal F,\PP)$. The covariates are a random variable $X\colon\Omega\to\X$, with $(\X,\A)$ a standard Borel space; this guarantees regular conditional distributions and allows mixed vectors of continuous, categorical and missing variables.

\begin{center}
\begin{tabular}{@{}ll@{}}
\toprule
Symbol & Meaning \\
\midrule
$D\in\{0,1\}$ & Channel serving the unit (1 = premium, 0 = basic) \\
$S\in\{0,1\}$ & Membership in the target group, with $\{S=1\}\subseteq\{D=1\}$ \\
$Y(0),\,Y(1)$ & Potential outcomes if served by the basic or the premium channel \\
$Y$ & Observed outcome \\
$P$ & Law of $X$ in the target group: $P(A)=\PP(X\in A\mid S=1)$ \\
$Q$ & Law of $X$ in the donor channel: $Q(A)=\PP(X\in A\mid D=0)$ \\
\bottomrule
\end{tabular}
\end{center}

We assume $\PP(S=1)>0$ and $\PP(D=0)>0$. A typical case is $S=\ind\{D=1,\ s(X)\le c\}$: the target group consists of the premium-channel units whose score $s(X)$ falls below a threshold, for instance the tail with the lowest purchase propensity.

\paragraph{Estimand.} The counterfactual mean of the target group in the basic channel,
\[
\theta=\E\bigl[Y(0)\mid S=1\bigr].
\]
The loss from diverting the group is $L=\E[Y\mid S=1]-\theta$; its first term is observed directly. When units generate several observations, $P$, $Q$ and $\theta$ are taken at the observation level (Section~\ref{sec:inference}).

\paragraph{Regressions.} Let $m_0$ be a version of $\E[Y\mid X,D=0]$, that is, an $\A$-measurable, $Q$-integrable function such that, for all $A\in\A$,
\[
\E\bigl[Y\,\ind\{X\in A\}\mid D=0\bigr]=\int_A m_0\dd Q .
\]
It exists by the Radon--Nikodym theorem applied to the measure $A\mapsto\E[Y\ind\{X\in A\}\mid D=0]$, which is absolutely continuous with respect to $Q$, and it is unique $Q$-almost surely. With discrete $X$, $m_0(x)$ is the rate of profile $x$; the definition through sets is needed because, with continuous $X$, each single value has probability zero. Similarly, $m_0^S$ is a version of $\E[Y(0)\mid X,S=1]$, unique $P$-almost surely.

\paragraph{Assumptions.}
\begin{enumerate}[label=(A\arabic*)]
  \item \emph{Consistency.} $Y=Y(0)$ on $\{D=0\}$. It fails if the basic channel changes when it receives the diverted units, for instance through congestion.
  \item \emph{Conditional mean exchangeability.} $m_0^S=m_0$ $P$-almost surely, where, by (A1), $m_0$ is also a version of $\E[Y(0)\mid X,D=0]$. It concerns only the basic channel: among units with the same profile, those in the target group would obtain there, on average, the same outcome as those it serves today. It says nothing about $Y(1)$ and cannot be tested with the data.
  \item \emph{Support.} $P\ll Q$: every set of profiles with probability zero in the donor channel has probability zero in the target group.
  \item \emph{Integrability.} $\E[|Y(0)|\mid S=1]<\infty$ and $\E[|Y|\mid D=0]<\infty$; it holds if $Y$ is bounded.
\end{enumerate}

\begin{remark}[A2 presupposes A3]\label{rem:a2a3}
Since $m_0$ is determined only $Q$-almost surely, the $P$-almost-sure equality in (A2) depends on the chosen version unless $P\ll Q$. If $P$ has mass where $Q$ has none, $m_0$ is arbitrary there and no donor data pin it down. This reappears in Theorem~\ref{thm:lebesgue}: the unsupported part of the target group is not estimable, not merely hard to estimate.
\end{remark}

\begin{remark}[Structural decomposition of the loss]\label{rem:decomposition}
If the outcome is a sum of components, $Y=\sum_kY_k$ (for example, purchases of different products), and the basic channel cannot produce the components in a set $K_{\mathrm{ex}}$, that is, $Y_k(0)=0$ almost surely for $k\in K_{\mathrm{ex}}$, then, by linearity,
\[
L=\sum_{k\in K_{\mathrm{ex}}}\E[Y_k\mid S=1]+\sum_{k\notin K_{\mathrm{ex}}}\bigl(\E[Y_k\mid S=1]-\E[Y_k(0)\mid S=1]\bigr).
\]
Reading $L$ as the causal loss requires consistency also in the premium channel ($Y=Y(1)$ on $\{D=1\}$). The first sum is observed without any identification assumption; only the second needs Sections~\ref{sec:identification} and~\ref{sec:classifier}. An error $e$ in $\theta_{\mathrm{com}}=\sum_{k\notin K_{\mathrm{ex}}}\E[Y_k(0)\mid S=1]$ produces an error $-e$ in $\hat L$, so the conclusion about $L$ is insensitive to failures of (A2) or (A3) when $\theta_{\mathrm{com}}/L$ is small. A large structural part is not enough on its own: with $a_{\mathrm{ex}}$ and $a_{\mathrm{com}}$ the two sums of $\E[Y_k\mid S=1]$ and $\rho=a_{\mathrm{ex}}/L$, we have $\theta_{\mathrm{com}}/L=a_{\mathrm{com}}/L-(1-\rho)$, which can be large even when $\rho$ is close to 1. Two caveats: if the unit later obtains the component through another route, $L$ is an upper bound on the loss to the organization; and with a binary ``buys anything'' outcome the decomposition is not additive, so one must work with the component indicators. The full swap, which also moves a group from the basic channel into the premium one, is the same construction with the roles of the channels exchanged.
\end{remark}
\section{Identification as a change of measure}\label{sec:identification}

Under (A1)--(A4), the counterfactual mean $\theta$ is the mean of the observed outcome in the donor channel weighted by $w=dP/dQ$. No part of the argument uses densities with respect to Lebesgue measure.

\begin{theorem}[Identification]\label{thm:identification}
Assume (A1)--(A4) and, in addition, $\int_\X\E\bigl[|Y|\mid X=x,D=0\bigr]\,P(dx)<\infty$. Let $w=dP/dQ$, which exists by (A3). Then
\[
\theta=\E\bigl[w(X)\,Y\mid D=0\bigr]=\frac{\E[\tilde w(X)\,Y\mid D=0]}{\E[\tilde w(X)\mid D=0]}\qquad\text{for every }\tilde w=c\,w,\ c>0.
\]
\end{theorem}
The additional condition holds whenever $Y$ is bounded.

\begin{proof}
\emph{Step 1 (tower).} Since the law of $X$ given $S=1$ is $P$, $\theta=\E[m_0^S(X)\mid S=1]=\int m_0^S\dd P$, finite by (A4) because $|m_0^S|\le\E[|Y(0)|\mid X,S=1]$ $P$-almost surely.
\emph{Step 2 (exchangeability).} By (A2), $\theta=\int m_0\dd P$.
\emph{Step 3 (Radon--Nikodym).} By (A3) there is a measurable $w\ge0$ with $P(A)=\int_A w\dd Q$ for all $A$. The change-of-measure formula, valid for indicators, for simple functions by linearity and for nonnegative functions by monotone convergence, extends to every $P$-integrable $g$: $\int g\dd P=\int g\,w\dd Q$. With $g=m_0$, $\theta=\int m_0\,w\dd Q$.
\emph{Step 4 (back to the data).} The additional condition is equivalent to $\E[w(X)|Y|\mid D=0]<\infty$, so $w(X)Y$ is integrable under $\PP(\cdot\mid D=0)$, and taking out what is known, $\E[w(X)Y\mid D=0]=\E[w(X)m_0(X)\mid D=0]=\int m_0\,w\dd Q=\theta$.
The self-normalized form follows from $\E[w(X)\mid D=0]=\int w\dd Q=P(\X)=1$: the ratio equals $\theta$ for $w$ and is invariant to multiplying $w$ by $c>0$.
\end{proof}

This invariance is what allows us, in Section~\ref{sec:classifier}, to use weights known only up to a constant.

\begin{example}\label{ex:toy}
Let three profiles A, B and C have $Q=(0.5,\,0.3,\,0.2)$, $P=(0.1,\,0.3,\,0.6)$ and $m_0=(4\,\%,\,2\,\%,\,0.5\,\%)$. The naive donor mean is $2.7\,\%$, but $\theta=0.1\cdot4\,\%+0.3\cdot2\,\%+0.6\cdot0.5\,\%=1.3\,\%$. The weights $w=P/Q=(0.2,\,1,\,3)$ recover it from the donor channel: $0.5\cdot0.2\cdot4\,\%+0.3\cdot1\cdot2\,\%+0.2\cdot3\cdot0.5\,\%=1.3\,\%$. The naive mean doubles $\theta$ because the donor channel consists mostly of A profiles, which respond more.
\end{example}

\begin{proposition}[Double robustness]\label{prop:dr}
Under (A1)--(A4), let $\bar w\ge0$ and $\bar m$ be measurable, with the following integrals finite, and
\[
\psi(\bar w,\bar m)=\E\bigl[\bar w(X)\,(Y-\bar m(X))\mid D=0\bigr]+\int\bar m\dd P .
\]
Then $\psi(\bar w,\bar m)-\theta=\int_\X(\bar w-w)(m_0-\bar m)\dd Q$. In particular, $\psi(\bar w,\bar m)=\theta$ if $\bar w=w$ or $\bar m=m_0$, $Q$-almost surely.
\end{proposition}
\begin{proof}
By Step 4 of the previous proof applied to $\bar w$, the first term equals $\int\bar w(m_0-\bar m)\dd Q$. By Step 3, $\int\bar m\dd P-\theta=\int(\bar m-m_0)\,w\dd Q$. Adding gives the result.
\end{proof}

The error is a product of the two estimation errors. This structure is what allows slow rates for each model and still $\sqrt n$ normality (Theorem~\ref{thm:normality}).
\section{The classifier identifies the Radon--Nikodym derivative}\label{sec:classifier}

A calibrated classifier that separates the target group from the donor channel yields $w$ up to a constant, and the set where that classifier is certain, $\{p=1\}$, is exactly the unsupported part of the target group.

\paragraph{Construction.} We pool both populations with a label: $(\tilde X,Z)$ with $Z\sim\mathrm{Bernoulli}(\pi)$, $\pi\in(0,1)$, $\tilde X\mid Z=1\sim P$ and $\tilde X\mid Z=0\sim Q$. The law of $\tilde X$ is the mixture $M=\pi P+(1-\pi)Q$, and the ideal classifier is a version $p(x)$ of $\PP(Z=1\mid\tilde X=x)$. In practice $\pi$ is the fraction of the target group in the pooled sample; it does not affect the estimand.

\begin{lemma}\label{lem:mixture}
$P\ll M$ and $Q\ll M$, with $dP/dM=p/\pi$ and $dQ/dM=(1-p)/(1-\pi)$, $M$-almost surely.
\end{lemma}
\begin{proof}
If $M(A)=0$, then $\pi P(A)+(1-\pi)Q(A)=0$ and, both coefficients being positive, $P(A)=Q(A)=0$. Moreover, for all $A\in\A$, $\pi P(A)=\PP(Z=1,\tilde X\in A)=\int_A p\dd M$ by definition of conditional expectation; hence $p/\pi$ is a version of $dP/dM$. The case of $Q$ is identical with $1-p$.
\end{proof}

\begin{theorem}[Odds ratio]\label{thm:odds}
Take a version of $p$ with values in $[0,1]$, and let $r=\frac{1-\pi}{\pi}\,\frac{p}{1-p}$ on $\{p<1\}$ and $r=0$ on $\{p=1\}$. Then $Q(\{p=1\})=0$ and, for all $A\in\A$, $\int_A r\dd Q=P(A\cap\{p<1\})$. In particular, under (A3), $r=dP/dQ$ $Q$-almost surely.
\end{theorem}
\begin{proof}
By Lemma~\ref{lem:mixture}, $Q(\{p=1\})=\int_{\{p=1\}}\frac{1-p}{1-\pi}\dd M=0$. For $A\in\A$, by the chain rule,
\[
\int_A r\dd Q=\int_{A\cap\{p<1\}}\frac{1-\pi}{\pi}\,\frac{p}{1-p}\cdot\frac{1-p}{1-\pi}\dd M=\int_{A\cap\{p<1\}}\frac{p}{\pi}\dd M=P(A\cap\{p<1\}).
\]
Under (A3), $P(\{p=1\})=0$ because $Q(\{p=1\})=0$, so $\int_A r\dd Q=P(A)$ for all $A$.
\end{proof}

Combined with the invariance in Theorem~\ref{thm:identification}, the odds $p/(1-p)$ suffice: the constant $(1-\pi)/\pi$ cancels. With a classifier of log-odds $F$, the weight is $e^{F}$. The theorem requires $p$ to be the true conditional probability: a classifier that ranks well but is miscalibrated orders units correctly and still gives biased weights.

\paragraph{Example~\ref{ex:toy} (continued).}
With $\pi=0.2$, the ideal classifier takes the values $p=(0.048,\,0.200,\,0.429)$ in A, B and C, with odds $(0.05,\,0.25,\,0.75)=w/4$, and the self-normalized estimator gives $1.3\,\%$. The classifier $\tilde p=p^2$ ranks identically (same AUC), but its odds $(0.0023,\,0.042,\,0.225)$ give $0.89\,\%$, a 32\,\% bias. Calibration matters, not only discrimination.


\begin{theorem}[Support and Lebesgue decomposition]\label{thm:lebesgue}
Without assuming (A3):
\begin{enumerate}
  \item $P\ll Q$ if and only if $P(\{p=1\})=0$.
  \item The Lebesgue decomposition of $P$ with respect to $Q$ is $P=\Pac+\Ps$, with $\Pac(A)=P(A\cap\{p<1\})$, $\Ps(A)=P(A\cap\{p=1\})$ and $d\Pac/dQ=r$.
  \item Let $\varepsilon=P(\{p=1\})$. If $\varepsilon<1$, (A2) holds on $\{p<1\}$ and $\int_{\{p<1\}}\E[|Y|\mid X=x,D=0]\,P(dx)<\infty$, self-normalized weighting with $r$ recovers
  \[
  \frac{\E[r(X)\,Y\mid D=0]}{\E[r(X)\mid D=0]}=\thac:=\E\bigl[Y(0)\mid S=1,\ p(X)<1\bigr],
  \]
  and $\theta=(1-\varepsilon)\thac+\varepsilon\,\ths$, with $\ths=\E[Y(0)\mid S=1,\ p(X)=1]$, defined when $\varepsilon>0$, not identified.
\end{enumerate}
\end{theorem}
\begin{proof}
(2) By Theorem~\ref{thm:odds}, $\Pac$ has density $r$ with respect to $Q$, hence $\Pac\ll Q$; $\Ps$ is concentrated on $\{p=1\}$, which is $Q$-null, hence $\Ps\perp Q$; by uniqueness of the Lebesgue decomposition, this is the decomposition. The set $\{p=1\}$ depends on the version of $p$ only up to $M$-null sets, which are null for both $P$ and $Q$. (1) $P\ll Q$ if and only if $\Ps=0$. (3) By Theorem~\ref{thm:odds}, $\E[r(X)\mid D=0]=1-\varepsilon$; repeating Steps 3 and 4 of Theorem~\ref{thm:identification} with $\Pac$, the numerator equals $\int_{\{p<1\}}m_0\dd P=(1-\varepsilon)\thac$. The relation with $\theta$ is the law of total probability.
\end{proof}

\begin{corollary}[Partial identification bounds]\label{cor:bounds}
Under the hypotheses of Theorem~\ref{thm:lebesgue}(3), if $Y\in[0,1]$, then $(1-\varepsilon)\thac\le\theta\le(1-\varepsilon)\thac+\varepsilon$, and the bias of weighting is $\thac-\theta=\varepsilon(\thac-\ths)$.
\end{corollary}
The width of the interval is $\varepsilon$, which can exceed the rate itself when rates are small: if an unsupported profile is added to Example~\ref{ex:toy} with $P=(0.1,\,0.3,\,0.5,\,0.1)$, then $\varepsilon=0.1$, $\thac=1.39\,\%$ and the interval is $[1.25\,\%,\,11.25\,\%]$. It is informative only under an additional assumption; for instance, if unsupported units do not respond better than the rest of the group, $\ths\le\thac$, the interval shrinks to $[(1-\varepsilon)\thac,\thac]$, here $[1.25\,\%,\,1.39\,\%]$, and the weighting estimator is an upper bound for $\theta$.

\begin{corollary}[Doubly robust estimator without support]\label{cor:drnosupport}
Under the hypotheses of Theorem~\ref{thm:lebesgue}, let $\bar m$ be $P$-integrable with $\bar m=m_0$ $P$-almost surely on $\{p<1\}$. With the weights $r$ (the last term is $0$ when $\varepsilon=0$),
\[
\frac{\E[r(X)(Y-\bar m(X))\mid D=0]}{\E[r(X)\mid D=0]}+\int\bar m\dd P=(1-\varepsilon)\thac+\varepsilon\,\E_P\bigl[\bar m(X)\mid p(X)=1\bigr].
\]
\end{corollary}
\begin{proof}
By Theorem~\ref{thm:odds}, the numerator of the first term is $\int_{\{p<1\}}(m_0-\bar m)\dd P=0$. The second term splits into $\{p<1\}$, where it integrates to $(1-\varepsilon)\thac$, and $\{p=1\}$, where it integrates to $\varepsilon\,\E_P[\bar m\mid p=1]$.
\end{proof}

Unlike the weighting estimator, which keeps the supported part, the doubly robust estimator fills the unsupported part with the extrapolation of the outcome model. It recovers $\theta$ if and only if that extrapolation is right on average, $\E_P[\bar m\mid p=1]=\ths$, an additional, untestable assumption. \citet{manamkoong2025} propose making it explicit under a smoothness assumption and bounding its effect with a sensitivity analysis.

\paragraph{Practical reading.} In finite samples one never observes $\hat p=1$ exactly. A natural diagnostic is the estimated mass $\hat P(\hat p>1-\delta)$ for small $\delta$; an AUC close to 1 is a warning sign, because it indicates that $P$ and $Q$ are close to mutually singular, not that the classifier is good.
\section{Target group defined by a known cutoff}\label{sec:cutoff}

In the typical case, the target group is a cutoff on a score the organization knows, $S=D\cdot\ind\{s(X)\le c\}$. Then no classifier separating the target group from the donor channel is needed: the ordinary channel propensity suffices, with the cutoff imposed by hand.

\begin{proposition}[Known-cutoff weights]\label{prop:cutoff}
Let $S=D\cdot\ind\{s(X)\le c\}$, with $s$ measurable and known, and let $e$ be a version of $\PP(D=1\mid X)$. Under (A3), there is $k>0$ such that
\[
\frac{dP}{dQ}(x)=k\,\ind\{s(x)\le c\}\,\frac{e(x)}{1-e(x)}\qquad Q\text{-almost surely.}
\]
\end{proposition}
\begin{proof}
Let $\mu$ be the law of $X$ in the population. Since $s(X)$ is a function of $X$, $\PP(S=1\mid X)=\ind\{s(X)\le c\}\,e(X)$. By the definition of conditional probability, $P(A)=\int_A\ind\{s\le c\}\,e\dd\mu/\PP(S=1)$ and $Q(A)=\int_A(1-e)\dd\mu/\PP(D=0)$ for all $A\in\A$, so $P$ and $Q$ are absolutely continuous with respect to $\mu$, with densities proportional to $\ind\{s\le c\}\,e$ and to $1-e$. The set $\{e=1\}$ is $Q$-null, and off it, under (A3), the chain rule gives $dP/dQ=(dP/d\mu)/(dQ/d\mu)=k\,\ind\{s\le c\}\,e/(1-e)$, with $k=\PP(D=0)/\PP(S=1)$.
\end{proof}

The identity is elementary, but it changes the estimation problem. The classifier of Theorem~\ref{thm:odds} must learn a sharp jump at $\{s=c\}$, with the few target-group cases as positives. The channel propensity is usually a smooth function and is estimated with all the units of the premium channel. The constant $k$ cancels under normalization, and the support diagnostic of Theorem~\ref{thm:lebesgue} is read from $\hat e$: support fails where $\hat e\to1$ inside $\{s\le c\}$. Lemma~\ref{lem:calibration} is the same identity applied to $(X,U)$.

\paragraph{Relation to the literature.} The identity is not new in substance. In subgroup analysis with propensity scores it is common practice to estimate the propensity in the full cohort and apply it within the subgroup \citep{rassen2012}, and \citet{dong2020} study when that model is preferable to one fitted within each subgroup, which balances better within the subgroup at the cost of more variance. What we add is the comparison with the ``target versus donors'' classifier, the usual choice when the change of measure is framed as density-ratio estimation \citep{bickel2009,sugiyama2012}, in the case where the target group is a cutoff on a score of the covariates themselves. With flexible models, both approaches are consistent and the differences are finite-sample ones, but they affect things the subgroup literature does not examine: the sign of trimming, the effective sample size and the validity of the bootstrap (Section~\ref{sec:res-cutoff}).
\section{Channel variables and differential measurement}\label{sec:channel}

A covariate recorded in only one channel breaks support even if the units are comparable. If it does not predict the outcome given the others, dropping it restores identification; if it confounds, neither using it nor dropping it identifies $\theta$.

\paragraph{Measurement model.} Each unit has latent covariates $X^*=(V,C^*)$, where $C^*$ takes values in $\mathcal C\cup\{\vacio\}$ and $C^*=\vacio$ means ``no actual data''. What is recorded depends on the channel: $X=(V,C)$, with $C=C^*$ if $D=1$ and $C=\vacio$ if $D=0$. More generally, $X=\varphi_D(X^*)$ with different measurement instruments, $\varphi_0\ne\varphi_1$. Examples: a commercial history kept only by premium-channel advisors, or diagnostic tests performed only in hospital. We write $P^*$ and $Q^*$ for the laws of $X^*$ in the target group and in the donor channel.

\begin{lemma}[Singular mass]\label{lem:singular}
$\varepsilon=\Ps(\X)=\sup\{P(B):\ B\in\A,\ Q(B)=0\}$.
\end{lemma}
\begin{proof}
Let $N$ be a $Q$-null set on which $\Ps$ is concentrated. For every $B$ with $Q(B)=0$, $\Pac(B)=0$ because $\Pac\ll Q$; hence $P(B)=\Ps(B)\le\Ps(\X)$. Moreover $P(N)=\Pac(N)+\Ps(N)=\Ps(\X)$, so the supremum is attained at $N$.
\end{proof}

\begin{proposition}[A channel variable breaks support]\label{prop:breaks}
Under the measurement model, let $\kappa=\PP(C^*\ne\vacio\mid S=1)$ be the fraction of the target group with actual data. Then $\varepsilon\ge\kappa$, and the ideal classifier equals $p=1$ $M$-almost surely on the profiles with $C\ne\vacio$. If $\kappa>0$, (A3) fails regardless of how the units compare.
\end{proposition}
\begin{proof}
Let $B=\mathcal V\times\mathcal C$, the set of profiles with recorded data. By the measurement model, $Q(B)=\PP(C\ne\vacio\mid D=0)=0$, whereas $P(B)=\kappa$; Lemma~\ref{lem:singular} gives $\varepsilon\ge\kappa$. Since $\PP(Z=0,\tilde X\in B)=(1-\pi)Q(B)=0$, $p=1$ $M$-almost surely on $B$.
\end{proof}

\begin{proposition}[Dropping a variable that does not predict the outcome]\label{prop:drop}
Assume:
\begin{enumerate}
  \item \emph{Latent support:} $P^*\ll Q^*$.
  \item \emph{Latent (A2):} $\E[Y(0)\mid X^*,S=1]=\E[Y(0)\mid X^*,D=0]$, $P^*$-almost surely.
  \item \emph{$C^*$ does not predict the outcome in the basic channel given $V$:} $\E[Y(0)\mid X^*,D=0]=\E[Y(0)\mid V,D=0]$, $Q^*$-almost surely.
\end{enumerate}
Then, with $V$ as the only covariates, (a) $P_V\ll Q_V$ and (b) $\E[Y(0)\mid V,S=1]=\E[Y(0)\mid V,D=0]$ $P_V$-almost surely, so Theorems~\ref{thm:identification}--\ref{thm:lebesgue} apply with $X=V$.
\end{proposition}
\begin{proof}
(a) If $Q_V(A)=0$, then $Q^*(A\times(\mathcal C\cup\{\vacio\}))=0$, and by (1) also $P^*(A\times(\mathcal C\cup\{\vacio\}))=P_V(A)=0$.
(b) Let $g$ be a version of $\E[Y(0)\mid V,D=0]$. By (3), $\E[Y(0)\mid X^*,D=0]=g(V)$ except on a $Q^*$-null set, which by (1) is also $P^*$-null. By (2), $\E[Y(0)\mid X^*,S=1]=g(V)$ $P^*$-almost surely, and since $V$ is a function of $X^*$, the tower property within $\{S=1\}$ gives $\E[Y(0)\mid V,S=1]=g(V)$.
\end{proof}

Condition (3) is an outcome-based covariate selection criterion \citep{deluna2011,vanderweele2011}: it suffices even if $C^*$ is associated with the channel.

\paragraph{When the variable does confound.} If (1) and (2) hold but (3) fails, let $\mu^*(x^*)=\E[Y(0)\mid X^*=x^*,D=0]$; then $\theta=\int\mu^*\dd P^*$. Using only $V$, the estimator converges to $\int g\dd P_V=\int g(v)\dd P^*$, with bias
\[
\int g\dd P_V-\theta=-\int_{\X^*}(\mu^*-g)\dd P^*.
\]
Under $Q^*$ this integral is zero by definition of $g$. Under $P^*$ it can be nonzero only if $C^*$ predicts $Y(0)$ given $V$ and its distribution given $V$ differs between groups: the two conditions of the classical definition of a confounder, necessary but not sufficient, because the effects could offset. The exact condition for (A2) with $V$ as covariates is $\E_{P^*}[\mu^*-g\mid V]=0$ $P_V$-almost surely; the bias itself vanishes if and only if $\int(\mu^*-g)\dd P^*=0$, which can also happen by cancellation. (3) is sufficient for both and does not involve $P^*$. When the bias is nonzero, neither option identifies $\theta$: without $C$, (A2) fails, and with $C$, (A3) fails and, as shown next, (A2) on the supported part generally fails too. Hence the bounds of Corollary~\ref{cor:bounds}, which assume (A2) there, do not apply either. What remains is a sensitivity analysis on (3).

\paragraph{Keeping the variable does not fix the supported part.} With $C$ among the covariates, the supported part is $\{C=\vacio\}$. But ``$\vacio$'' does not mean the same in both groups: in the donor channel it mixes units with $C^*=\vacio$ and $C^*\ne\vacio$, because $C$ is never recorded there; in the target group it means $C^*=\vacio$. Since the whole donor channel has $C=\vacio$, $m_0(v,\vacio)=g(v)$, and the weighting estimator converges to $\E_{P^*}[g(V)\mid C^*=\vacio]$, which in general is neither $\theta$ nor $\thac$: (A2) on $\{C=\vacio\}$ requires $\mu^*(V,\vacio)=g(V)$ $P$-almost surely, that is, that among donors with the same $V$ those with and without a history respond alike in the basic channel; (3) implies it, but it is weaker. The doubly robust estimator, in addition, extrapolates the outcome model to units with data (Corollary~\ref{cor:drnosupport}). That model is fitted in the donor channel, where $C$ never varies; if it does not depend on $C$ outside the support, as with trees, which cannot split on a variable that is constant in training, it extrapolates $g(V)$, and the doubly robust estimator converges to $\E_{P^*}[g(V)]$, as if $C$ were not used.

\paragraph{Testability.} Condition (3) concerns $Y(0)$ in the basic channel, where $C^*$ is not observed, so it cannot be tested directly. That $C$ predicts the outcome in the premium channel is a hint but not proof, because $Y(1)$ is observed there.

\paragraph{Relation to the missingness-pattern approach.} \citet{dagostino2000} and \citet{blake2020} estimate the propensity within each missingness pattern and require the partially observed variable not to confound in the pattern where it is missing. A channel variable is the case in which the treatment determines the pattern: in the observed pattern there are no donor units (Proposition~\ref{prop:breaks}), and in the missing one it requires $\mu^*(V,\vacio)=g(V)$, which (3) implies.

\paragraph{Practical diagnostic.} Removing a channel variable should visibly reduce the classifier's AUC and the estimated mass $\hat P(\hat p>1-\delta)$; a large jump in that mass when a variable is included is the signature of Proposition~\ref{prop:breaks}.

\subsection{Different measurement instruments}

A variable missing in one channel is the most visible case of a general problem: each channel measures with its own instrument. Let $X^*$ be the latent covariate, taking values in a standard Borel space $\X^*$, with laws $P^*$ and $Q^*$ in the target group and the donor channel. Each channel records $X=\varphi_D(X^*)$, with $\varphi_0,\varphi_1\colon\X^*\to\X$ measurable, so that $P=P^*\circ\varphi_1^{-1}$ and $Q=Q^*\circ\varphi_0^{-1}$. Let $\mu^*(x^*)=\E[Y(0)\mid X^*=x^*,D=0]$.

\begin{theorem}[Different measurement instruments]\label{thm:measurement}
Assume (A1), latent support $P^*\ll Q^*$ and latent (A2): $\E[Y(0)\mid X^*,S=1]=\mu^*(X^*)$ $P^*$-almost surely.
\begin{enumerate}
  \item \emph{Common variables.} For every measurable $\psi$, $P^*\circ\psi^{-1}\ll Q^*\circ\psi^{-1}$. In particular, if $\psi\circ\varphi_0=\psi\circ\varphi_1=k$ (variables both channels measure the same way), the laws of $K=\psi(X)$ in the two groups satisfy $P_K\ll Q_K$.
  \item \emph{Identification with the common variables.} If, in addition, $\mu^*(X^*)=\E[Y(0)\mid K,D=0]$ $Q^*$-almost surely, Theorems~\ref{thm:identification}--\ref{thm:lebesgue} apply with $K$ as covariates.
  \item \emph{Differential measurement with support.} If $\varphi_0$ is injective and $P\ll Q$ (on standard Borel spaces, $\varphi_0$ then has a Borel image and a measurable inverse by the Lusin--Souslin theorem, and $\varphi_1(X^*)$ lies in that image $P^*$-almost surely because $Q$ is concentrated on it), the weighting estimator with $X$ converges, for bounded $Y$, to
  \[
  \theta_X=\int_{\X^*}\mu^*\bigl(\varphi_0^{-1}(\varphi_1(x^*))\bigr)\,P^*(dx^*),\qquad
  \theta_X-\theta=\int_{\X^*}\bigl[\mu^*\circ\varphi_0^{-1}\circ\varphi_1-\mu^*\bigr]\dd P^*.
  \]
  \item \emph{Indistinguishability.} Under the conditions of 3, there is a world without measurement error, with $\varphi_1=\varphi_0$ and the latent law of the target group equal to $P^*\circ(\varphi_0^{-1}\circ\varphi_1)^{-1}$, that reproduces the law of the target-group and donor-channel data (and of all observed data if, in addition, the recorded covariates of every premium unit have a law absolutely continuous with respect to $Q$), satisfies the maintained assumptions and has $\theta=\theta_X$. No procedure based only on the observed data distinguishes $\theta$ from $\theta_X$.
\end{enumerate}
\end{theorem}
\begin{proof}
(1) If $Q^*(\psi^{-1}(B))=0$, latent support gives $P^*(\psi^{-1}(B))=0$. If $\psi\circ\varphi_0=\psi\circ\varphi_1=k$, then $K=k(X^*)$ in both channels, and it suffices to apply the above with $\psi=k$.

(2) This is Proposition~\ref{prop:drop} with $V=K$, a function of $X^*$.

(3) In the donor channel, $X=\varphi_0(X^*)$ and $\varphi_0$ is injective with a measurable inverse, so conditioning on $X$ is equivalent to conditioning on $X^*$; by (A1), $m_0(x)=\mu^*(\varphi_0^{-1}(x))$ $Q$-almost surely, and since $P\ll Q$ this version is fixed $P$-almost surely. The weighting estimator converges to $\int m_0\dd P$ (Theorem~\ref{thm:identification}, which uses only $P\ll Q$ for this convergence), and by the change of variables for image measures, $\int m_0\dd P=\int m_0(\varphi_1(x^*))\,P^*(dx^*)=\theta_X$. By latent (A2) and the tower property, $\theta=\int\mu^*\dd P^*$.

(4) In the alternative world the donor channel does not change. The target group is recorded with $\varphi_0$ and its latent law is $\tilde P^*=P^*\circ(\varphi_0^{-1}\circ\varphi_1)^{-1}$, so its observed law is $P^*\circ(\varphi_0\circ\varphi_0^{-1}\circ\varphi_1)^{-1}=P^*\circ\varphi_1^{-1}=P$, the same as in the real world. For its observed outcome it suffices to take, given $X$, the same conditional law as in the real world, and under latent (A2) its response in the basic channel is $\mu^*$; its $\theta$ is $\int\mu^*\dd\tilde P^*=\theta_X$. The alternative world satisfies the assumptions: since $\varphi_0$ is injective, $Q^*(A)=Q(\varphi_0(A))$ and $\tilde P^*(A)=P(\varphi_0(A))$, so $P\ll Q$ gives $\tilde P^*\ll Q^*$, and latent (A2) holds by construction. Under the additional condition, premium units outside the target group are treated in the same way, transforming their latent law by $\varphi_0^{-1}\circ\varphi_1$, which is then defined almost surely; without it, $\varphi_1(X^*)$ may fall outside the image of $\varphi_0$ for those units, and a world with $\varphi_1=\varphi_0$ cannot reproduce them.
\end{proof}

\paragraph{Reading.} A variable measured differently is more dangerous than a missing one. The missing one breaks support and leaves a trace: the classifier is certain, $p=1$, and $\hat\varepsilon>0$. The differently measured one can preserve support, and the classifier interprets the measurement difference as a real difference between units. In an example with $P^*=N(1,1)$, $Q^*=N(0,1)$, $\mu^*(x)=\sigma(-3-0.8x)$ and a premium channel that records the variable shifted by $c$, a shift of $c=\pm0.3$ standard deviations produces biases of $-20.5\,\%$ and $+25.4\,\%$ on $\theta=2.92\,\%$. With a logistic classifier, the estimator matches $\theta_X$ ($z=+0.6$ in both cases). The AUC changes (0.76 without error; 0.82 and 0.69 with it), but in a way that a real population difference would also produce. The only way out is information external to the design: units measured with both instruments, for instance customers who switched channels, which allow estimating $\varphi_0^{-1}\circ\varphi_1$.
\section{Estimation and inference}\label{sec:inference}

\paragraph{Sampling.} We observe an i.i.d. pooled sample $O_i=(Z_i,X_i,(1-Z_i)Y_i)$, $i=1,\dots,n$, with $\PP(Z=1)=\pi$, $X\mid Z=1\sim P$ and $X\mid Z=0\sim Q$; for $\theta$ only the outcomes of donor units are needed. When a unit (a customer) generates several observations (calls), $O_i$ is the whole block of that unit and everything below applies at the unit level. We assume that the channel and target-group membership are properties of the unit, constant across its observations, and that $P$, $Q$, $m_0$ and $\theta$ are defined at the observation level, so that each unit weighs according to its number of observations. If that number depends on $Y(0)$ given $X$, the unit-level and observation-level estimands, and the corresponding versions of (A2), differ.

\paragraph{Estimators.} With $\hat o=\hat p/(1-\hat p)$ and $n_1$ the number of target-group observations,
\[
\hat\theta_{\mathrm{H}}=\frac{\sum_{i:Z_i=0}\hat o(X_i)\,Y_i}{\sum_{i:Z_i=0}\hat o(X_i)},\qquad
\hat\theta_{\mathrm{DR}}=\frac{\sum_{i:Z_i=0}\hat o(X_i)\,(Y_i-\hat m(X_i))}{\sum_{i:Z_i=0}\hat o(X_i)}+\frac1{n_1}\sum_{i:Z_i=1}\hat m(X_i).
\]
The first is the Hájek weighting estimator and the second the sample version of $\psi$ in Proposition~\ref{prop:dr} with the weighting term normalized. Both use cross-fitting: units are split into $K$ folds and, for observations in fold $k$, $\hat p$ and $\hat m$ are trained on the other folds. With a known cutoff, $\hat o$ is replaced by $\ind\{s\le c\}\,\hat e/(1-\hat e)$ (Proposition~\ref{prop:cutoff}), and the sample also includes the premium units outside the target group. Under population sampling of $(D,S,X,Y)$, the corresponding influence function is $\PP(S=1)^{-1}\bigl[S(m_0-\theta)+(1-D)\,\ind\{s\le c\}\,e/(1-e)\,(Y-m_0)\bigr]$; we expect an analogue of Theorem~\ref{thm:normality} but do not prove it.

\subsection{Weight trimming}\label{sec:trimming}

Trimming the weights changes the estimand. For $\tau>0$, let $w_\tau=\min(w,\tau)$ and $\varphi_\tau(t)=\min(1,\tau/t)$ for $t>0$. With the true weights, the trimmed self-normalized estimator converges to
\[
\theta_\tau=\frac{\E[w_\tau(X)\,Y\mid D=0]}{\E[w_\tau(X)\mid D=0]}.
\]

\begin{proposition}[Trimming bias]\label{prop:trimming}
Under (A1)--(A4):
\begin{enumerate}
  \item $\displaystyle\theta_\tau-\theta=\frac{\Cov_P\bigl(m_0(X),\varphi_\tau(w(X))\bigr)}{\E_P[\varphi_\tau(w(X))]}$.
  \item If $g(t)=\E_P[m_0(X)\mid w(X)=t]$ has a nonincreasing version, $\theta_\tau\ge\theta$; if it has a nondecreasing version, $\theta_\tau\le\theta$.
  \item If $Y\ge0$, the unnormalized version is always below: $\E[w_\tau(X)\,Y\mid D=0]\le\theta$.
\end{enumerate}
\end{proposition}
\begin{proof}
Since $P(\{w=0\})=\int_{\{w=0\}}w\dd Q=0$, on $\{w>0\}$ we have $w_\tau=\varphi_\tau(w)\,w$. By the tower property and the change of measure, $\E[w_\tau(X)Y\mid D=0]=\int w_\tau m_0\dd Q=\E_P[\varphi_\tau(w)\,m_0]$ and $\E[w_\tau(X)\mid D=0]=\E_P[\varphi_\tau(w)]\in(0,1]$; subtracting $\theta=\E_P[m_0]$ gives (1). For (2), the tower property gives $\Cov_P(m_0,\varphi_\tau(w))=\Cov_P(g(T),\varphi_\tau(T))$ with $T=w(X)$. The function $\varphi_\tau$ is nonincreasing; if $g$ is too, $(g(T)-g(T'))(\varphi_\tau(T)-\varphi_\tau(T'))\ge0$ for an independent copy $T'$, and its expectation is $2\Cov_P(g(T),\varphi_\tau(T))$ (Chebyshev's association inequality); the nondecreasing case is symmetric. For (3), $\varphi_\tau\le1$ and $m_0\ge0$ give $\E_P[\varphi_\tau(w)m_0]\le\theta$.
\end{proof}

\paragraph{Interpretation.} A large $w(x)$ means that profile $x$ is much more frequent in the target group than in the donor channel. There is no default direction: it depends on whether, within the target group, the profiles most typical of the premium channel respond more or less in the basic channel. If they respond less, $g$ is decreasing and trimming overestimates $\theta$ and hence underestimates the loss; in our simulation they respond more (Section~\ref{sec:res-trimming}). The condition concerns $g$ as a function of the weight, not of $X$, and it can be checked approximately, because $m_0$ is identified on the support.

\paragraph{Example~\ref{ex:toy} (trimming).} With $\tau=2$, the weights $(0.2,\,1,\,3)$ become $(0.2,\,1,\,2)$ and $\theta_\tau=1.2\,\%/0.8=1.5\,\%$, against $\theta=1.3\,\%$. The formula agrees: $\varphi_\tau=(1,\,1,\,2/3)$, $\E_P[\varphi_\tau]=0.8$ and the covariance is $0.16$ points, so the bias is $0.2$ points. The unnormalized version gives $1.2\,\%$, below $\theta$.

\begin{corollary}[Trimming with the doubly robust estimator]\label{cor:drtrimming}
By Proposition~\ref{prop:dr} with $\bar w=w_\tau$, $\psi(w_\tau,\bar m)-\theta=-\int_{\{w>\tau\}}(w-\tau)(m_0-\bar m)\dd Q$.
\end{corollary}
The corollary concerns the unnormalized form $\psi$: trimming biases only to the extent that the outcome model is wrong in the trimmed region. For the normalized form we use, the bias is $\E_Q[w_\tau(m_0-\bar m)]/\E_Q[w_\tau]-\E_P[m_0-\bar m]$, which vanishes if $\bar m=m_0$ on the whole support, but not necessarily if the model is right only in the trimmed region. Since Proposition~\ref{prop:dr} holds for any weights, when estimated weights are trimmed the bias depends on both errors, and it vanishes if $\bar m=m_0$.

\paragraph{Estimated weights.} In practice $\hat w$ is trimmed, and trimming removes both true weight and estimation noise, so the sign in Proposition~\ref{prop:trimming} is not guaranteed. In the companion note, included in the repository that accompanies this paper, we prove that it is preserved if the estimation error is a profile-independent multiplicative noise with a log-concave density. Before trimming, the effect of trimming computed with $\hat m$ instead of $Y$, or $\hat m$ by mass deciles of $\hat w$, indicates which way it will push. \citet{crump2009} trim on the propensity and redefine the estimand; \citet{ma2020} keep the estimand and characterize the bias of trimming by discarding in the unnormalized estimator, of which item 3 is the version for capping; the contribution of Proposition~\ref{prop:trimming} is the normalized case, with its exact formula. \citet{liu2023} treat truncation as a change of estimand for the effect on the treated.

\subsection{Variance and asymptotic normality}

\paragraph{Influence function.} Let $o=p/(1-p)$ and $\sigma^2(x)=\Var(Y\mid X=x,D=0)$. The influence function of the doubly robust estimator, of the same form as the efficient one for the counterfactual mean of the treated \citep{hahn1998}, is
\begin{gather*}
\phi(O)=\frac1\pi\Bigl[Z\,(m_0(X)-\theta)+(1-Z)\,o(X)\,(Y-m_0(X))\Bigr],\\
V=\E[\phi^2]=\frac{\Var_P(m_0(X))}{\pi}+\frac{\E_Q[w(X)^2\sigma^2(X)]}{1-\pi},
\end{gather*}
where the cross term vanishes because $Z(1-Z)=0$ and the second term uses $o=\pi w/(1-\pi)$ (Theorem~\ref{thm:odds}). This closed form is for i.i.d. observations; with several observations per unit, see the appendix.

\begin{proposition}[When the variance is finite]\label{prop:variance}
If $Y$ is bounded and $0<\sigma_-^2\le\sigma^2\le\sigma_+^2$ $Q$-almost surely, then $V<\infty$ if and only if $\E_Q[w^2]<\infty$, that is, if and only if $\chi^2(P\|Q)=\E_Q[w^2]-1<\infty$. Absolute continuity $P\ll Q$ does not imply this condition.
\end{proposition}
\begin{proof}
The first term of $V$ is finite because $m_0$ is bounded; the second lies between $\sigma_-^2\,\E_Q[w^2]/(1-\pi)$ and $\sigma_+^2\,\E_Q[w^2]/(1-\pi)$. To see that $P\ll Q$ is not enough, let $Q$ be uniform on $(0,1)$ and $w(x)=x^{-1/2}/2$, which integrates to 1; then $P\ll Q$ but $\E_Q[w^2]=\int_0^1(4x)^{-1}dx=\infty$.
\end{proof}

Thus (A3) suffices to identify $\theta$ but not to estimate it at rate $n^{-1/2}$: with the conditional variance bounded below, one also needs $\chi^2(P\|Q)<\infty$, and standard theory requires something stronger, bounded weights.

\paragraph{Effective sample size.} Kish's $\neff=(\sum_i\hat w_i)^2/\sum_i\hat w_i^2$ over the $n_0$ donor observations satisfies, with the true weights, $\neff/n_0\to1/\E_Q[w^2]=1/(1+\chi^2(P\|Q))$, so it estimates the $\chi^2$ divergence; with estimated weights it can move either way (noise in $\hat w$ lowers it, shrinkage by a regularized classifier raises it), and in our simulation it falls below the theoretical value except under low overlap, where the theoretical value is itself imprecise (Section~\ref{sec:design}). In finite samples it is always finite, even when $\chi^2=\infty$, and it then goes to zero very slowly and noisily: in the counterexample of Proposition~\ref{prop:variance}, the median over replications goes from 0.45 with $n_0=10^3$ to about 0.22 with $n_0=10^7$, and with $n_0=10^3$ the 1st and 99th percentiles across replications are 0.05 and 0.65. A reasonable $\neff$ does not guarantee that $\chi^2$ is finite; one should also examine the stability of the upper quantiles of $\hat w$.

\begin{theorem}[Asymptotic normality of the doubly robust estimator]\label{thm:normality}
Assume (A1)--(A4), bounded $Y$ and, in addition:
\begin{enumerate}
  \item \emph{Strong overlap:} $p(X)\le1-\eta$ $Q$-almost surely for some $\eta>0$, that is, bounded $w$.
  \item Cross-fitting by unit with a fixed number $K$ of folds of size proportional to $n$, with $\hat p\le1-\eta$ and $\hat m$ bounded by the same bounds as $Y$.
  \item \emph{Rates:} $r_o=\max_k\|\hat o_k-o\|$ and $r_m=\max_k\|\hat m_k-m_0\|$ in $L^2(Q)$ go to 0 in probability, and $r_o\,r_m=o_P(n^{-1/2})$.
\end{enumerate}
Then $\sqrt n(\hat\theta_{\mathrm{DR}}-\theta)\to N(0,V)$ in distribution, and the sample mean of $\hat\phi(O_i)^2$ estimates $V$ consistently.
\end{theorem}
The proof, in Appendix~\ref{app:proof}, adapts Theorem 3.1 of \citet{chernozhukov2018} and its example of the effect on the treated: by Proposition~\ref{prop:dr}, the bias is bilinear in the errors of the auxiliary parameters (Neyman orthogonality), cross-fitting controls the empirical-process term without complexity conditions, and the denominator is handled by the delta method.

\paragraph{Beyond strong overlap.} Condition 1 is stronger than (A3) and than $\chi^2<\infty$. If the weights are unbounded but $\chi^2<\infty$, $V$ is finite, although the standard proof does not cover the case. If $\chi^2=\infty$ and the conditional variance is bounded below, no regular estimator converges at rate $\sqrt n$: this is the irregular identification case of \citet{khan2010}. \citet{ma2020} show that the limit may then be non-normal and propose subsampling inference, and \citet{ma2023} give adaptive inference for doubly robust estimators with trimming and bias correction. Our simulation lies in this regime: the channel index is unbounded inside the target region, so the weights are unbounded while $\chi^2$ is finite, and the implementation clips $\hat p$ to $[10^{-6},1-10^{-6}]$ rather than to $1-\eta$; its coverage results are evidence outside the conditions of Theorem~\ref{thm:normality}.

\paragraph{Intervals.} We use $\hat\theta_{\mathrm{DR}}\pm1.96\,(\hat V/n)^{1/2}$. Theorem~\ref{thm:normality} covers neither weight trimming nor the bootstrap, which we evaluate in the simulation (Sections~\ref{sec:res-trimming} and~\ref{sec:res-bootstrap}). If the target group is defined with an estimated threshold $\hat c$, its uncertainty adds a term to $V$ that we do not include: we treat the selection rule as fixed.
\section{Sensitivity to unobserved confounding}\label{sec:sensitivity}

If (A2) fails, $\theta$ is no longer identified. The marginal sensitivity model bounds how far $\theta$ can move from the identified value as a function of a single parameter $\Lambda\ge1$.

\paragraph{Model.} We follow \citet{tan2006} and \citet{dorn2023}; \citet{zhao2019} give sensitivity intervals for weighting estimators under this model via the percentile bootstrap. Suppose there is an unobserved variable $U$ such that exchangeability holds given $(X,U)$: $\E[Y(0)\mid X,U,S=1]=\E[Y(0)\mid X,U,D=0]$. In the pooled sample of Section~\ref{sec:classifier}, let $\lambda(X,U)$ be the ratio of the odds of belonging to the target group given $(X,U)$ and given $X$ alone. The model MSM($\Lambda$) assumes $\Lambda^{-1}\le\lambda(X,U)\le\Lambda$: knowing $U$ can multiply or divide those odds by at most $\Lambda$. With $\Lambda=1$ we recover (A2).

\begin{proposition}[Sensitivity bounds]\label{prop:sensitivity}
Assume (A1), (A3), (A4), $Y\ge0$ and MSM($\Lambda$), and let $\theta_0=\int m_0\dd P$ be the value identified under (A2).
\begin{enumerate}
  \item \emph{Validity.} $\theta\in[\theta_0/\Lambda,\ \Lambda\theta_0]$.
  \item \emph{Sharpness for rare outcomes.} If, in addition, $\PP(Y>0\mid X,D=0)\le1/(1+\Lambda)$ $Q$-almost surely, both bounds are attained: there is a world compatible with MSM($\Lambda$) and with everything observed in which $\theta=\Lambda\theta_0$, and another in which $\theta=\theta_0/\Lambda$.
  \item \emph{General binary case.} For $Y\in\{0,1\}$ with $m=m_0(x)$, the sharp bounds integrate with respect to $P$ the values $\min\bigl(\Lambda m,\,1-(1-m)/\Lambda\bigr)$ (upper) and $\max\bigl(m/\Lambda,\,1-\Lambda(1-m)\bigr)$ (lower); this is the quantile rule of \citet{dorn2023} carried over to $\theta$.
\end{enumerate}
\end{proposition}
\begin{proof}
By Bayes' rule, $\lambda(x,\cdot)$ is the density of the law of $U$ given $X=x$ in the target group with respect to that in the donor channel, so $\E_Q[\lambda\mid X]=1$, and the derivative of the law of $(X,U)$ in the target group with respect to that in the donor channel is $w(X)\,\lambda(X,U)$. Applying Theorem~\ref{thm:identification} to $(X,U)$, $\theta=\E_Q[w(X)\,\E_Q(\lambda Y\mid X)]$.

(1) Since $Y\ge0$ and $\Lambda^{-1}\le\lambda\le\Lambda$, $m_0(X)/\Lambda\le\E_Q[\lambda Y\mid X]\le\Lambda m_0(X)$; integrate against $w\dd Q$.

(2) Let $\pi_+(X)=\PP(Y>0\mid X,D=0)$. For the upper bound, take $U=Y(0)$, $\lambda=\Lambda$ if $Y(0)>0$ and $\lambda=c(X)=(1-\Lambda\pi_+)/(1-\pi_+)$ if $Y(0)=0$. Then $\E_Q[\lambda\mid X]=1$; moreover $c\le1\le\Lambda$, and $c\ge1/\Lambda$ is equivalent to $\pi_+\le1/(1+\Lambda)$, so MSM($\Lambda$) holds. Since $\E_Q[\lambda\mid X]=1$, neither the law of $X$ in the target group nor the donor data change, and the target group's $Y(0)$ is never observed, so the world is compatible with the observables; in it, $\E_Q[\lambda Y\mid X]=\Lambda m_0(X)$ and $\theta=\Lambda\theta_0$. For the lower bound, $\lambda=1/\Lambda$ if $Y(0)>0$ and $\lambda=(1-\pi_+/\Lambda)/(1-\pi_+)$ if $Y(0)=0$; this value is $\ge1$ and is $\le\Lambda$ if $\pi_+\le\Lambda/(1+\Lambda)$, which holds because $\pi_+\le1/(1+\Lambda)\le\Lambda/(1+\Lambda)$.

(3) Let $\bar\lambda_1$ and $\bar\lambda_0$ be the means of $\lambda$ over $\{Y=1\}$ and $\{Y=0\}$ given $X$. The constraint $\E_Q[\lambda\mid X]=1$ reads $m\bar\lambda_1+(1-m)\bar\lambda_0=1$ with $\bar\lambda_0,\bar\lambda_1\in[1/\Lambda,\Lambda]$, and $\E_Q[\lambda Y\mid X]=m\bar\lambda_1$. Solving for $\bar\lambda_0$, the maximum of $m\bar\lambda_1$ is $\min(\Lambda m,\,1-(1-m)/\Lambda)$ and the minimum is $\max(m/\Lambda,\,1-\Lambda(1-m))$; both are attained with $\lambda$ constant on each class. We also checked the formula against the linear-programming solution in 2{,}000 random cases, with maximum error $10^{-15}$.
\end{proof}

\paragraph{Practical use.} Given a confidence interval $[\hat L,\hat U]$ for $\theta_0$, the interval $[\hat L/\Lambda,\ \hat U\Lambda]$ combines sampling and confounding uncertainty. The \emph{breakdown $\Lambda$} is the smallest $\Lambda$ that changes the conclusion; for instance, the loss $\E[Y\mid S=1]-\theta$ remains positive for every $\Lambda<\E[Y\mid S=1]/\theta_0$.

\paragraph{Calibrating $\Lambda$.} Comparing with observed covariates, dropping each from the propensity model and measuring how much the odds change, is established practice \citep{mcclean2024,baitairian2025}. When the target group is a cutoff on a score of the covariates themselves, this must be done with the channel model, not with the ``target versus donors'' model.

\begin{lemma}[Calibration with the channel model]\label{lem:calibration}
Let $S=D\cdot\ind\{s(X)\le c\}$ with $s$ observed. On the region $R=\{s(X)\le c\}$,
\[
\lambda(X,U)=\frac{\odds(D=1\mid X,U)}{\odds(D=1\mid X)}.
\]
\end{lemma}
\begin{proof}
In the pooled sample, by Bayes' rule, $\odds(Z=1\mid x,u)=\frac{\pi}{1-\pi}\cdot\frac{\PP(S=1\mid x,u)/\PP(S=1)}{\PP(D=0\mid x,u)/\PP(D=0)}$, and likewise without $u$. On $R$, $\PP(S=1\mid x,u)=\PP(D=1\mid x,u)$ and $\PP(S=1\mid x)=\PP(D=1\mid x)$; the constants cancel in the ratio.
\end{proof}

Outside $R$, the target-group model is dominated by the cutoff, which is not confounding. Inside $R$, the lemma allows calibrating with an ordinary channel propensity model. We define the \emph{benchmark} of covariate $j$ as $\Lambda_j(q)$, the $q$-quantile on $R$ of $\max(\lambda_j,1/\lambda_j)$, with $\lambda_j=\odds(D=1\mid X)/\odds(D=1\mid X\text{ without }j)$, both estimated out of sample. We use quantiles rather than the sample maximum, which grows with the sample size; they are a scale for comparison, not a verification of the model. Two caveats: that a hidden confounder is weaker than the strongest observed covariate is an assumption, not a consequence of the data; and dropping covariates one at a time understates confounding when several correlated covariates share it \citep{mcclean2024}.
\section{Simulation design}\label{sec:design}

The simulation uses a generator with known $\theta$ and varies one factor at a time around a base scenario, so that each theoretical result has a scenario that tests it against a prior numerical prediction.

\paragraph{Generator.}
\begin{enumerate}
  \item \emph{Covariates:} $V=(V_1,\dots,V_4)$ normal with correlations $0.3^{|j-k|}$, plus a three-level categorical $V_5$ whose distribution depends on $V_1$; all are customer characteristics. An unobserved $U$, independent of $V$, enters only the unobserved-confounding scenario.
  \item \emph{Latent channel variable:} $H=\ind\{C^*\ne\vacio\}$ indicates an actual history, with $\PP(H=1\mid V)=\sigma(-0.4+0.8V_1+0.3V_3)$; if $H=1$, the history takes the value $N(0.5V_1,1)$. It is recorded if $D=1$ and $H=1$, and $C=\vacio$ otherwise.
  \item \emph{Channel:} $\PP(D=1\mid V,H,U)=\sigma(\alpha_0+\lambda\,\alpha^\top V+\zeta H+\nu U)$. The parameter $\lambda$ controls overlap, and $\zeta\ne0$ makes the distribution of $C^*$ given $V$ differ between channels, a necessary condition for $C^*$ to confound.
  \item \emph{Outcomes:} two common products, $\PP(Y_k(d)=1\mid X^*,U)=\sigma(\beta_{k0}+\beta_{kd}+\beta_k^\top V+\delta H+\nu'U+e)$ for $k=1,2$, with $e$ a customer random effect, and a product exclusive to the premium channel with $Y_3(0)=0$. The outcome used by the estimators is $Y=Y_1+Y_2$.
  \item \emph{Target group:} $S=\ind\{D=1,\ s^*(V)\le c^*\}$, with $s^*$ the true purchase-propensity index in the premium channel and $c^*$ its 10th percentile in that channel. The rule is fixed, so $\theta$ does not change across replications.
  \item \emph{Units and observations:} each customer generates between 1 and 5 calls, a number drawn independently of all other variables, with channel and covariates constant across them; the unit- and observation-level estimands therefore coincide.
  \item \emph{Truth:} $\theta$ and the theoretical limit of each estimator are computed on a population of 1.5 million customers; the true conditional means and $dP_V/dQ_V$ are obtained by integrating out $H$, $U$ and the customer effect with Gauss--Hermite quadrature. Their Monte Carlo error, about 0.001 points (0.006 under confounding through the history and 0.015 under unobserved confounding), is negligible relative to that of the replication means. Trimming effects, much smaller, are predicted as paired differences within the population.
\end{enumerate}

\paragraph{Scenarios and predictions.} The base scenario excludes $C$, has $\zeta\ne0$ and $\delta=\nu=\nu'=0$, high overlap ($\lambda=0.5$; $\chi^2=3.6$), LightGBM with five-fold cross-fitting by customer for $\hat p$ and $\hat m$, and 30{,}000 customers per replication: about 62{,}000 donor calls and 2{,}800 target-group calls. Around it we vary the inclusion of $C$, confounding through the history ($\delta=1.2$), overlap ($\lambda=1.5$ and $3$; $\chi^2=6.3$ and about 36; the latter is unstable because $w^2$ is very heavy-tailed: on six populations of 1.5 million customers it ranged from 33 to 63, so the theoretical $\neff/n_0$ under low overlap, 0.03, is only indicative), model specification, unobserved confounding ($\nu=\nu'=1$), trimming at the 99th, 95th and 90th percentiles of the donor weights, and size (3{,}000 to 30{,}000 customers). Each scenario has 200 replications, so that 95\,\% coverage is estimated with a standard error of about 1.5 points. Predictions that do not follow directly from the theory were registered, with a date, before running the corresponding scenario; the code, the predictions (time-stamped in the version-control history of the repository) and the raw results will be made public with the paper.

\paragraph{The AUC does not measure overlap.} When calibrating $\lambda$, the AUC of the ideal classifier stayed between 0.89 and 0.93 for $\lambda$ between 0.3 and 2, while $\chi^2$ ranged from 3.5 to about 11, not monotonically (it is smallest for $\lambda$ between 0.5 and 1). The definition of the target group, a cutoff on a score, already separates it from the donor channel. We therefore set the levels by $\chi^2$, which is what governs the variance (Proposition~\ref{prop:variance}).
\section{Simulation results}\label{sec:results}

\subsection{Identification, channel variables and specification}\label{sec:res-main}

In the eight scenarios of Table~\ref{tab:main}, the means of the estimators converge to the predicted limit. With small samples, dispersion and coverage fail, not the mean (Section~\ref{sec:res-bootstrap}). A discrepancy under medium overlap ($z=-3.0$ for Hájek and $-2.1$ for DR) did not reappear in a replication with 200 new seeds ($z=+0.8$ and $+1.5$; with all 400 replications, $-1.4$ and $-0.3$). This replication was not planned in advance; with about 40 comparisons, one $|z|\ge3$ by chance has a probability of roughly 10\,\%.

\begin{table}[ht]
\centering\small
\caption{Main results. Rates in \%; in parentheses, $z=$ (mean $-$ prediction)/(Monte Carlo error). SD: standard deviation of the DR estimator across replications. Coverage of the DR Wald interval at 95\,\%, with respect to $\theta$. In the scenarios with $C$, the theoretical $\neff/n_0$ is computed with the supported part; it is practically the same as without $C$ (0.217), because $\PP(H=0\mid V,D=1)$ varies little on the target region.}\label{tab:main}
\resizebox{\textwidth}{!}{%
\begin{tabular}{@{}lcccccc c@{}}
\toprule
Scenario & $\theta$ & Prediction & Hájek & DR & SD & Coverage & $\neff/n_0$ (theoretical) \\
\midrule
Base & 0.958 & 0.958 & 0.958 (0.0) & 0.961 (0.4) & 0.112 & 95.5\,\% & 0.14 (0.22) \\
With $C$ & 0.958 & H 0.940; DR 0.958 & 0.939 ($-$0.1) & 0.958 (0.0) & 0.107 & 96.5\,\% & 0.12 (0.22) \\
$C$ confounds, without $C$ & 1.712 & 1.314 & 1.313 ($-$0.1) & 1.320 (0.6) & 0.132 & 15.5\,\% & 0.14 (0.22) \\
$C$ confounds, with $C$ & 1.712 & H 1.281; DR 1.314 & 1.280 (0.0) & 1.320 (0.6) & 0.129 & 17.5\,\% & 0.12 (0.22) \\
Medium overlap & 1.178 & 1.178 & 1.151 ($-$3.0) & 1.158 ($-$2.1) & 0.136 & 94.5\,\% & 0.10 (0.14) \\
Low overlap & 1.274 & 1.274 & 1.283 (0.5) & 1.298 (1.1) & 0.300 & 91.0\,\% & 0.03 (0.03) \\
Unobserved confounding & 2.697 & 1.193 & 1.198 (0.6) & 1.204 (1.3) & 0.127 & 0.0\,\% & 0.14 (0.20) \\
3{,}000 customers & 0.958 & 0.958 & 0.887 ($-$1.2) & 0.883 ($-$1.2) & 0.912 & 70.0\,\% & 0.02 (0.22) \\
\bottomrule
\end{tabular}}
\end{table}

\begin{itemize}
  \item \emph{Identification.} Without confounding or a channel variable, Hájek and DR recover $\theta$ and the interval covers 95.5\,\%.
  \item \emph{Channel variable.} With $C$, the estimated unsupported mass was $\hat\varepsilon=0.326$, against $\kappa=0.335$ (Proposition~\ref{prop:breaks}), and Hájek converges to $\thac$ (the weighting limit of Section~\ref{sec:channel} coincides with $\thac$ here because $\delta=0$), with the predicted bias of $-0.018$ points. The original prediction also assigned $\thac$ to the DR estimator and failed ($z=2.4$ and $4.3$). The DR estimator extrapolates the outcome model to the unsupported part (Corollary~\ref{cor:drnosupport}); the prediction derived from this (0.958 and 1.314) was added after the ``With $C$'' scenario and registered before running ``$C$ confounds, with $C$'', where it held ($z=0.6$).
  \item \emph{Confounding through the history.} Without $C$, the bias is as predicted: $-0.40$ points, 23\,\% of $\theta$. With $C$, $\theta$ is not recovered either, as Section~\ref{sec:channel} says.
  \item \emph{Unobserved confounding.} The estimators converge to the limit predicted from the observables, less than half of $\theta$, with zero coverage. The standard error is well calibrated (97\,\% of the actual dispersion) and the diagnostics resemble those of the base scenario: nothing in the data reveals the problem.
  \item \emph{Overlap.} With $\chi^2=36$, the standard deviation almost triples, the estimated standard error is 87\,\% of the actual dispersion and coverage drops to 91\,\%.
\end{itemize}

\begin{table}[ht]
\centering\small
\caption{Specification (base scenario). Rates in \%.}\label{tab:specification}
\begin{tabular}{@{}llllcc@{}}
\toprule
Classifier & Outcome model & Estimator & Prediction & Mean ($z$) & Coverage \\
\midrule
LightGBM & LightGBM & DR & 0.958 & 0.961 (0.4) & 95.5\,\% \\
LightGBM & linear in $V_1$ & DR & 0.958 & 0.955 ($-$0.3) & 95.0\,\% \\
Linear logistic & LightGBM & DR & 0.958 & 0.946 ($-$2.3) & 88.0\,\% \\
Linear logistic & --- & Hájek & 1.024 & 1.023 ($-$0.4) & --- \\
Linear logistic & linear in $V_1$ & DR & 1.157 & 1.150 ($-$1.5) & 13.5\,\% \\
--- & linear in $V_1$ & Regression & 1.069 & 1.075 (1.4) & --- \\
\bottomrule
\end{tabular}
\end{table}

With the flexible classifier and the misspecified (linear) outcome model, the DR estimator is unbiased and covers 95\,\%; with both misspecified, it converges to the predicted limit, far from $\theta$ (Table~\ref{tab:specification}). With the misspecified (linear) classifier and the flexible outcome model, a bias of $-0.012$ points and 88\,\% coverage remain: the outcome model is a LightGBM, not $m_0$, and the product of errors is not negligible at these sizes. Even so, the DR estimator reduces the bias of the corresponding Hájek estimator by 82\,\%.

\subsection{Trimming}\label{sec:res-trimming}

We measure the effect of trimming as the paired difference, trimmed minus untrimmed, within each replication, in thousandths of a percentage point. With the true weights and the true response $m_0$, in 60 replications of the base scenario, Proposition~\ref{prop:trimming} predicts $-0.37$, $-3.18$ and $-7.98$ at the 99th, 95th and 90th percentiles, and we observe $-0.37\pm0.02$, $-3.13\pm0.06$ and $-7.81\pm0.09$: the profiles with the largest weights respond more, and trimming lowers the estimate. With the weights estimated by the classifier, the sign reverses: $+10.9\pm0.7$ at the 90th percentile, positive in all 60 replications (with the simulated purchases, noise hides it: $+11.8\pm7.3$ in the same 60 replications). Within the DR estimator, the mean effect of trimming lies between $-0.1$ and $+2.1$ thousandths in the base and low-overlap scenarios, and between $+6.4$ and $+12.5$ under medium overlap, indistinguishable from zero ($|z|<2$); at the 90th percentile, the root mean squared error drops from 0.300 to 0.163 points under low overlap and from 0.111 to 0.082 in the base scenario. But the standard error computed with the trimmed weights is only 64\,\% of the actual dispersion under low overlap, and coverage falls to 78.5\,\%; with the untrimmed standard error it is 96.5--99.5\,\% at the 95th and 90th percentiles in the base and low-overlap scenarios, but only 78--88.5\,\% with 3{,}000 customers. The mechanism of the reversal is analyzed in the companion note.

\subsection{Known-cutoff weights}\label{sec:res-cutoff}

We compare the estimator with the weights of Proposition~\ref{prop:cutoff} with the standard one in the same replications: same data, same LightGBM and same outcome model; only the weights change. The channel propensity is estimated adding the premium calls outside the target group, which the standard estimator does not use. The predictions were registered before running.

\begin{table}[ht]
\centering\small
\caption{Known-cutoff weights in the base scenario, by size. 200 replications; $\theta=0.958\,\%$. SD in points. The last column is the effect of trimming the Hájek estimator at the 90th percentile computed with $m_0$, in thousandths, over these 200 replications (the $+10.9$ in the text comes from a separate set of 60); with the true weights it is $-7.8$.}\label{tab:cutoff}
\begin{tabular}{@{}llccccc@{}}
\toprule
Customers & Weights & DR (\%) & SD & Coverage & $\neff/n_0$ & Trimming p90 \\
\midrule
30{,}000 & standard & 0.961 & 0.112 & 95.5\,\% & 0.14 & $+10.4$ \\
30{,}000 & known cutoff & 0.960 & 0.098 & 92.5\,\% & 0.20 & $-6.2$ \\
10{,}000 & standard & 0.962 & 0.294 & 88.0\,\% & 0.07 & $+19.8$ \\
10{,}000 & known cutoff & 0.957 & 0.195 & 91.0\,\% & 0.17 & $-4.8$ \\
3{,}000 & standard & 0.883 & 0.912 & 70.0\,\% & 0.02 & $+50.0$ \\
3{,}000 & known cutoff & 0.944 & 0.371 & 87.0\,\% & 0.11 & $-10.2$ \\
\bottomrule
\end{tabular}
\end{table}

Five of the six predictions hold (Table~\ref{tab:cutoff}): no appreciable bias; $\neff/n_0$ approaches the theoretical 0.22; the dispersion falls by 12, 34 and 59\,\%; trimming recovers the sign of the true weights; and coverage improves with 10{,}000 and 3{,}000 customers. Coverage with 30{,}000 fails: we predicted between 93 and 97\,\% and obtained $92.5\,\%\pm1.9$, because the standard error is 93\,\% of the actual dispersion. With 10{,}000 customers, the standard deviation is 15\,\% above what the $\sqrt n$ rate predicts from the same estimator with 30{,}000 ($0.098\cdot\sqrt3=0.170$), against 52\,\% for the standard one.

A subsequent analysis, not registered in advance, separates the two novelties: imposing the cutoff and using the premium units outside the target group. Imposing the cutoff on the standard classifier, without new data, does not reduce the dispersion (0.117, 0.298 and 0.933 points) and only partly corrects the sign of trimming ($+1.0$, $+2.7$ and $+9.6$). The improvement comes mainly from estimating the channel propensity with many more positive cases; the known cutoff is what makes using it possible.

\begin{table}[ht]
\centering\small
\caption{Known-cutoff weights in the other scenarios (200 replications, same data as Table~\ref{tab:main}). In each cell, standard $\to$ known cutoff. Trimming of the Hájek estimator at the 90th percentile computed with $m_0$, in thousandths.}\label{tab:cutoff-scenarios}
\begin{tabular}{@{}lccc@{}}
\toprule
Scenario & SD of DR & Coverage & Trimming p90 \\
\midrule
With $C$ & 0.107 $\to$ 0.096 & 96.5 $\to$ 93.5\,\% & $+16.4\to-0.4$ \\
$C$ confounds, without $C$ & 0.132 $\to$ 0.115 & 15.5 $\to$ 9.0\,\% & $+20.1\to-7.4$ \\
$C$ confounds, with $C$ & 0.129 $\to$ 0.113 & 17.5 $\to$ 6.5\,\% & $+32.0\to+3.9$ \\
Medium overlap & 0.136 $\to$ 0.147 & 94.5 $\to$ 93.5\,\% & $+19.6\to-15.6$ \\
Low overlap & 0.300 $\to$ 0.342 & 91.0 $\to$ 87.0\,\% & $+34.5\to-20.6$ \\
Unobserved confounding & 0.127 $\to$ 0.109 & 0.0 $\to$ 0.0\,\% & $+23.1\to+4.6$ \\
\bottomrule
\end{tabular}
\end{table}

In the other scenarios, with a second registered prediction (Table~\ref{tab:cutoff-scenarios}), the estimators converge to the same limits as the standard one ($|z|\le1.7$), because only the estimation of the same weights changes; with confounding, coverage with respect to $\theta$ is even lower, because the estimator is more precise around a wrong value. With the channel variable, the weights do not blow up: donors never have it recorded, and the lack of support shows up as $\hat e>0.99$ in 33.3\,\% of the target group, against $\kappa=0.335$. Under low overlap, trimming with these weights reproduces what Proposition~\ref{prop:trimming} predicts for the true weights ($-2.2$, $-12.1$ and $-20.6$ against $-3.0$, $-11.9$ and $-20.8$). The effect of trimming shrinks substantially in every scenario and becomes negative or about zero, except in ``$C$ confounds, with $C$'' and ``Unobserved confounding'', where it remains positive.

Two predictions fail: the dispersion does not decrease under medium overlap, and under low overlap it increases and coverage falls from 91 to 87\,\%. Our interpretation, untested: as overlap worsens, the channel propensity approaches 1 inside the target group and its odds are estimated from few donors, whereas the standard classifier, with very few positive cases, shrinks them toward moderate values. Trimming within the DR estimator equalizes both methods: under low overlap, the root mean squared error at the 90th percentile is 0.160 with the known cutoff and 0.163 with the standard estimator, with 97.5 and 98\,\% coverage using the untrimmed standard error. The precision gain is therefore limited to good overlap, especially with small samples.

\subsection{Sensitivity}\label{sec:res-sensitivity}

The smallest $\Lambda$ whose point interval contains $\theta$ is $\theta/\theta_0$: 1.30 when the history confounds and 2.26 with unobserved confounding. Coverage of the combined interval $[\hat L/\Lambda,\hat U\Lambda]$ rises quickly around that value (Table~\ref{tab:sensitivity}); it stays below the normal approximation because the standard error is smaller in replications with low estimates (correlations of 0.64 and 0.70). The rare-outcome condition of Proposition~\ref{prop:sensitivity}, checked through the bound $\PP(Y>0\mid V,D=0)\le g(V)=\E[Y\mid V,D=0]$ at the 99.99th percentile of $g$ in the donor channel (not at the essential supremum), holds at least up to $\Lambda=5.6$ in the base scenario, 1.8 under confounding through the history and 7.3 under unobserved confounding. With normal $U$, $\lambda$ is unbounded and no finite $\Lambda$ satisfies the MSM in the strict sense, so in that scenario Proposition~\ref{prop:sensitivity} guarantees nothing; $\Lambda=2.26$ contains $\theta$ by construction, and we use it only as a scale.

\begin{table}[ht]
\centering\small
\caption{Coverage of the combined sensitivity interval.}\label{tab:sensitivity}
\begin{tabular}{@{}lccccc@{}}
\toprule
Scenario & $\theta/\theta_0$ & $\Lambda=1.25$ & $\Lambda=1.5$ & $\Lambda=2$ & $\Lambda=2.5$ \\
\midrule
$C$ confounds, without $C$ & 1.30 & 87\,\% & 100\,\% & 100\,\% & 100\,\% \\
Unobserved confounding & 2.26 & 0\,\% & 3\,\% & 71\,\% & 99\,\% \\
\bottomrule
\end{tabular}
\end{table}

To calibrate $\Lambda$, in the unobserved-confounding scenario $U$ enters the channel with coefficient 1 and the strongest observed covariate, $V_1$, with 0.5. We predicted, before computing it, that the method of Lemma~\ref{lem:calibration} would recover the exact value for $U$ and that $U$ would exceed every observed covariate; both hold (Table~\ref{tab:calibration}). The $\Lambda$ needed to cover the truth, 2.26 (known here only because $\theta$ is), lies between the median and the 90th percentile of $V_1$: a confounder somewhat stronger than the strongest observed covariate suffices to explain the whole gap. This is not the breakdown $\Lambda$ of Section~\ref{sec:sensitivity}, which depends on the conclusion of interest and is computed from the data. In a second estimate with another sample, the ``target versus donors'' model agrees with the channel model at the median (2.06 against 1.99) but overstates the extremes: 6.34 and 29.98 at the 90th and 99th percentiles, against 5.37 and 15.51 with the channel model and 5.38 and 14.07 exact.

\begin{table}[ht]
\centering\small
\caption{Calibration of $\Lambda$ with the channel model: quantiles on $R$ of $\max(\lambda,1/\lambda)$.}\label{tab:calibration}
\begin{tabular}{@{}lccc@{}}
\toprule
Variable & Median & 90th percentile & 99th percentile \\
\midrule
$U$, exact & 1.99 & 5.38 & 14.07 \\
$U$, estimated & 2.03 & 5.21 & 14.14 \\
$V_1$ & 1.72 & 3.07 & 5.81 \\
$V_2$ & 1.17 & 1.53 & 2.17 \\
$V_3$ & 1.14 & 1.45 & 1.98 \\
$V_4$ & 1.12 & 1.36 & 1.70 \\
$V_5$ & 1.08 & 1.21 & 1.42 \\
\bottomrule
\end{tabular}
\end{table}

\subsection{Sample size and intervals: Wald and bootstrap}\label{sec:res-bootstrap}

With the standard estimator, the standard deviation of the DR estimator is 0.912, 0.508, 0.294, 0.183 and 0.112 points with 3{,}000, 5{,}000, 10{,}000, 15{,}000 and 30{,}000 customers: it falls like $n^{-0.91}$, much faster than $n^{-1/2}$, and Wald coverage is 70, 80, 88, 92 and 95.5\,\%. The predictions for 5{,}000 and 15{,}000 were registered in advance and held. The estimated $\neff/n_0$ follows the same progression (0.02, 0.07 and 0.14 with 3{,}000, 10{,}000 and 30{,}000, against the theoretical 0.22). This suggests a diagnostic, not yet tested: if $\neff/n_0$ keeps increasing with the sample size, the intervals are not yet reliable.

\begin{table}[ht]
\centering\small
\caption{Customer-level bootstrap: mean standard error in points and, in parentheses, 95\,\% coverage, on the same 40 replications.}\label{tab:bootstrap}
\begin{tabular}{@{}lcccc@{}}
\toprule
Customers & Actual SD & Wald & Bootstrap with replacement & Bayesian bootstrap \\
\midrule
30{,}000 & 0.109 & 0.102 (97.5\,\%) & 0.162 (100\,\%) & 0.147 (100\,\%) \\
10{,}000 & 0.274 & 0.230 (82.5\,\%) & 0.432 (100\,\%) & 0.370 (97.5\,\%) \\
\bottomrule
\end{tabular}
\end{table}

We evaluated the customer-level bootstrap with replacement, re-estimating LightGBM and the cross-fitting in each of 60 resamples, in 80 replications with 30{,}000 and 10{,}000 customers, and, in 40 of them, a Bayesian variant that gives each customer an $\mathrm{Exp}(1)$ weight instead of copies. All four registered predictions failed (Table~\ref{tab:bootstrap}). We predicted that the bootstrap with replacement would be close to the actual dispersion with 30{,}000 customers and between Wald and the actual dispersion with 10{,}000; both overstate it (over the 80 replications, 0.165 against 0.119 and 0.446 against 0.322). We attributed the excess to duplicated customers and predicted that the Bayesian variant would remove it, but it reduces the standard error by only 9 and 14\,\%, and the excess over the actual dispersion by 28 and 39\,\%.

A second check, also registered, points to the cause. Applied to the estimator with known-cutoff weights, whose weights are much less noisy, the Bayesian bootstrap is nearly calibrated: mean standard error of 0.101 against an actual dispersion of 0.096 on the same 40 replications with 30{,}000 customers, and 0.208 against 0.187 with 10{,}000, with coverage of 95 and 90\,\% ($\pm5$). This indicates that the excess is tied to the noise of the estimated weights: resampling reduces the effective information, and since the dispersion of the standard classifier grows quickly as information shrinks, the bootstrap world exaggerates it. This is evidence from a single data-generating process, not a proof.
\section{Illustrations with public data}\label{sec:illustrations}

We complement the simulation with two public data sets in which the truth is known. In the first, both versions of a differentially measured covariate exist for the same people, so the bias predicted by Theorem~\ref{thm:measurement} can be checked directly. In the second, an experimental control group provides the benchmark for $\theta$. No proprietary data are used.

\subsection{Self-reported versus measured body mass index}\label{sec:nhanes}

\paragraph{Data.} We use the 2017--2018 cycle of the US National Health and Nutrition Examination Survey \citep{nchs2020}: the 4{,}034 adults aged 20 or more with measured and self-reported height and weight, systolic blood pressure (mean of up to three readings), waist circumference, education and income-to-poverty ratio. We do not use the survey weights and treat the sample as the population. Self-reported BMI is computed from the reported weight and height and rounded to 0.1, the precision of the published measured BMI, for comparability. Self-reported BMI is on average 0.81 lower than measured BMI (standard deviation 1.99), and the gap grows with BMI: from 0.00 in the lowest quartile of measured BMI to $-1.66$ in the highest, in line with the literature on self-reported anthropometry \citep{connorgorber2007}. Using it in one channel only is differential measurement in the sense of Section~\ref{sec:channel}; that the error is not classical (it grows with BMI) is what makes the bias sizeable.

\paragraph{Design.} Each replication assigns every person to the target channel with probability $\sigma\bigl(-1+\beta_b z(\mathrm{BMI})+0.5\,z(\mathrm{age})-0.4\,z(\mathrm{income})\bigr)$, where BMI is the measured value, $z$ standardizes and $\beta_b=0.8$; about 30\,\% end up in the target group. The channel is synthetic, but covariates, outcomes and measurement error are real. Since the assignment depends only on the true covariates, latent exchangeability and support hold by construction, and $\theta$, the mean outcome of the target group, is observed in each replication. We compare four specifications, all with age, sex, race, education and income: (A) measured BMI in both groups; (B) self-reported BMI in the target group and measured BMI in the donors, the channel-measurement case of Theorem~\ref{thm:measurement}; (C) without BMI; (D) self-reported BMI in both groups. The estimator is the doubly robust one with five-fold cross-fitting, with a penalized logistic classifier and a ridge regression, both additive in cubic splines; there are 200 replications. The replications share the population and only the assignment varies, so we evaluate bias, not coverage.

\paragraph{Registered predictions.} Before running, we computed the limit of specification B from Theorem~\ref{thm:measurement}(3), with $\mu(V,b)=\E[Y\mid V,\mathrm{BMI}=b]$ fitted on the whole sample: averaging over the assignment probabilities $\pi_i$, the predicted bias is $\sum_i\pi_i[\mu(V_i,b_i^{\mathrm{self}})-\mu(V_i,b_i)]/\sum_i\pi_i$, namely $-0.47$\,mmHg for blood pressure and $-2.44$\,cm for the waist. Strictly, Theorem~\ref{thm:measurement} assumes deterministic instruments, and the self-report is not a function of the measured value; but its proof goes through unchanged here, because the assignment depends only on $(V,\mathrm{BMI})$, so $\theta=\int m_0(V,\mathrm{BMI})\dd P^*$ while the estimator converges to $\int m_0(V,\mathrm{BMI}^{\mathrm{self}})\dd P^*$. The registered criterion was $|z|<3$, allowing a model error of up to 25\,\% of the predicted bias. We also predicted no bias in A, no mass of $\hat p$ close to 1 in A or B, an AUC difference between A and B below 0.02, and in C a negative bias larger than that of B for the waist.

\begin{table}[ht]
\centering
\caption{NHANES, base design (200 replications). Mean bias of the DR estimator with its Monte Carlo error in parentheses. No replication had any target unit with $\hat p>0.99$.}\label{tab:nhanes}
\begin{tabular}{@{}lcccc@{}}
\toprule
Specification & Blood pressure (mmHg) & Waist (cm) & AUC & $\neff/n_0$ \\
\midrule
A: measured in both & $-0.11$ (0.06) & $-0.18$ (0.02) & 0.743 & 0.41 \\
B: differential & $-0.48$ (0.06) & $-2.56$ (0.02) & 0.715 & 0.50 \\
\quad prediction (Theorem~\ref{thm:measurement}) & $-0.47$ & $-2.44$ & & \\
C: without BMI & $-1.78$ (0.05) & $-10.70$ (0.04) & 0.653 & 0.74 \\
D: self-reported in both & $-0.47$ (0.06) & $-1.33$ (0.03) & 0.737 & 0.44 \\
\bottomrule
\end{tabular}
\end{table}

\paragraph{Results.} Table~\ref{tab:nhanes} confirms the prediction of Theorem~\ref{thm:measurement} approximately: $-0.48$ against $-0.47$\,mmHg ($z=-0.3$) and $-2.56$ against $-2.44$\,cm. The waist deviation, 5\,\%, is within the registered tolerance but statistically distinguishable from zero ($z=-5.7$); net of the small bias of A, $-2.37$ against $-2.44$. Nothing in the observed data signals the problem. No target unit has $\hat p>0.99$, and $\neff/n_0$ \emph{improves} from 0.41 to 0.50: the target group, selected for high BMI, reports a BMI shifted down toward the donors, so the groups separate less and the weights look better precisely when the estimator is biased (noise alone, in D, moves $\neff/n_0$ only to 0.44). The AUC falls from 0.743 to 0.715, but specification A is not available in practice, so this fall cannot be observed. Dropping BMI is much worse ($-10.70$\,cm), because BMI drives the assignment. Two registered predictions fail: A has a bias of $-0.18$\,cm in the waist ($z=-8.3$), which grows with the strength of the assignment and vanishes with LightGBM ($-0.01$; Table~\ref{tab:nhanes-robust}), so it comes from the spline models; and the AUC falls by more than the registered 0.02.

Specification D, registered as exploratory, suggests a practical lesson with limits: using the self-reported BMI in \emph{both} groups halves the bias of the waist relative to B, but not that of blood pressure ($-0.47$ against $-0.48$). Measuring both groups with the same imprecise instrument leaves residual confounding, but does not evaluate the outcome model at a shifted argument; which effect dominates depends on the outcome. Table~\ref{tab:nhanes-robust} shows that for the waist D beat B in all four designs, by 24 to 73\,\%, whereas for blood pressure it was worse in two of them. We have no general result on when this ordering holds.

\begin{table}[ht]
\centering
\caption{NHANES, robustness (200 replications each). Mean bias of the DR estimator in the waist (cm) and, in the last two columns, in blood pressure (BP); in parentheses, the registered prediction.}\label{tab:nhanes-robust}
\resizebox{\textwidth}{!}{%
\begin{tabular}{@{}lcccccc@{}}
\toprule
Design & A & B & C & D & B, BP (mmHg) & D, BP (mmHg) \\
\midrule
Weak assignment ($\beta_b=0.4$) & $-0.08$ & $-2.26$ ($-2.19$) & $-5.69$ & $-0.61$ & $-0.42$ ($-0.41$) & $-0.27$ \\
Base ($\beta_b=0.8$), splines & $-0.18$ & $-2.56$ ($-2.44$) & $-10.70$ & $-1.33$ & $-0.48$ ($-0.47$) & $-0.47$ \\
Strong assignment ($\beta_b=1.2$) & $-0.46$ & $-2.89$ ($-2.61$) & $-14.44$ & $-2.20$ & $-0.56$ ($-0.51$) & $-0.69$ \\
Base, LightGBM & $-0.01$ & $-2.42$ ($-2.44$) & $-10.38$ & $-1.23$ & $-0.20$ ($-0.47$) & $-0.45$ \\
\bottomrule
\end{tabular}}
\end{table}

\paragraph{Robustness.} We varied the strength with which BMI drives the assignment and replaced the splines by LightGBM, again registering the predictions before running (Table~\ref{tab:nhanes-robust}). For each design we registered five: B within the tolerance above; the bias of A in the waist below 20\,\% of that of B (a relative criterion, chosen after the base run had shown a model bias in A; under $|z|<3$, A would also fail in the weak and strong designs); no target unit with $\hat p>0.99$ in A or B; C negative and larger than B in the waist; and D smaller than B in the waist. Of these 15 predictions, 14 hold. In all four designs, no target unit has $\hat p>0.99$ and $\neff/n_0$ is larger in B than in A. The failure is blood pressure with LightGBM: $-0.20$ against $-0.47$ ($z=4.0$). Blood pressure depends weakly on BMI and is noisy, so LightGBM flattens that dependence and transmits only part of the measurement error, while the prediction used $\mu$ fitted with splines. An unregistered check supports this: with $\mu$ fitted by LightGBM, the predicted bias is $-0.31$, and the observed $-0.20$ is at $z=1.7$ from it (for the waist, $-2.40$ instead of $-2.44$). When the outcome depends weakly on the mismeasured covariate, the prediction itself depends on the outcome model.

\subsection{An experimental benchmark: NSW and CPS}\label{sec:lalonde}

\paragraph{Data.} We use the subsample of \citet{dehejia1999} of the National Supported Work experiment, with 185 participants and 260 randomized controls, and the CPS-1 comparison group of \citet{lalonde1986}, with 15{,}992 people. The target group is the participants; the donors, the CPS. The estimand is the mean of their 1978 earnings without the program, whose experimental benchmark is the mean of the controls: 4{,}555 dollars (standard error 340). The covariates are age, education, race, marital status, lack of a degree, earnings in 1974 and 1975 with indicators of zero earnings, and the squares of age and education. This analysis was not registered in advance; we report the median over five cross-fitting partitions.

\paragraph{Results.} The unadjusted CPS mean is 14{,}847 dollars. Weighting alone gives 5{,}092 with a logistic classifier and 3{,}990 with LightGBM, whereas the doubly robust estimator gives 4{,}783 and 4{,}716 (standard errors 397 and 876), both within one standard error of the benchmark (340). Overlap is poor: with LightGBM, $\neff$ is 33 out of 15{,}992 donors.

\paragraph{Differential measurement.} The two sources record earnings differently \citep{smith2005}. One difference is visible: CPS earnings are top-coded (19\,\% of 1974 earnings sit exactly at the cap of 25{,}862 dollars) and NSW earnings are not. Two participants (1.1\,\%) earned more than the cap in 1974, so their profiles have no support by construction: a singular mass in the sense of Lemma~\ref{lem:singular}, where the ideal classifier equals 1. Neither estimated classifier comes close: they give these two units $\hat p$ below 0.3, because the 19\,\% of donors at the cap make high earnings look donor-like. A support failure at the edge of a top-coded variable thus points the diagnostic the wrong way. Applying the cap to both groups changes the estimates by less than 30 dollars, so here it does not matter. The outcome is also top-coded: 26\,\% of CPS 1978 earnings sit at the cap, which biases the donor outcome model downward; applying the cap to the controls lowers the benchmark by 54 dollars, to 4{,}501, small relative to the standard errors. Other differences in how earnings were collected, if they preserve support, are not detectable at all (Theorem~\ref{thm:measurement}).
\section{Discussion}\label{sec:discussion}

\paragraph{Practical recommendations.} The theory and the simulation suggest six guidelines for evaluating a channel switch with observational data:
\begin{enumerate}
  \item If the target group is a cutoff on a known score, estimate the weights with the channel propensity, using all units of the premium channel, and impose the cutoff (Proposition~\ref{prop:cutoff}).
  \item Exclude from the covariates the variables recorded by only one channel when they do not predict the outcome given the others (Proposition~\ref{prop:drop}), and diagnose lack of support through the mass of $\hat p$ or $\hat e$ close to 1 (Proposition~\ref{prop:breaks}). When differential measurement preserves support, no check on the data detects it (Theorem~\ref{thm:measurement}), and the overlap diagnostics may even improve (Section~\ref{sec:nhanes}); external information is needed. If one channel measures a covariate less accurately and the accurate version cannot be obtained, consider using the less accurate version in both channels as well, and compare: in our illustration it reduced the bias for the waist but not for blood pressure, and we have no general result.
  \item Use the doubly robust estimator with cross-fitting by unit, bearing in mind that without support it extrapolates the outcome model (Corollary~\ref{cor:drnosupport}).
  \item For the interval, use Wald with large samples. With medium samples and known-cutoff weights, both Wald and the unit-level Bayesian bootstrap were close to nominal in the simulation (90--95\,\%, the bootstrap on only 40 replications); with the standard classifier, the bootstrap is conservative. Be wary if $\neff/n_0$ is far below what is expected.
  \item Do not carry the sign of trimming from the true weights over to the estimated ones without checking it. If trimming, do it within the doubly robust estimator and compute the standard error with the untrimmed weights.
  \item Accompany the estimate with a sensitivity analysis, with $\Lambda$ calibrated against the observed covariates through the channel model (Lemma~\ref{lem:calibration}).
\end{enumerate}

\paragraph{Limitations.} The simulation evidence comes from a single data-generating process, built so that each theoretical result had a scenario to test it, and the public-data illustrations are either semi-synthetic (the channel is assigned by us) or rest on a single experimental benchmark; quantitative conclusions from the simulation, such as the gains of the known cutoff or the excess of the bootstrap, should be read as illustrations of the mechanisms, not as general magnitudes. The sensitivity bounds assume the MSM, which does not hold in the strict sense with unbounded confounders. We treat the threshold that defines the target group as fixed. And the estimand is the counterfactual mean under the basic channel as it is today: if receiving the diverted units congests it, (A1) fails.

\paragraph{Open problems.} Four questions remain open: a theory explaining the excess of the bootstrap with noisy estimated weights and the validity of the Bayesian bootstrap with known-cutoff weights; the variance contribution of a threshold estimated as a quantile of the score, which relates to choosing the threshold itself; combining the sensitivity analysis with lack of support (Corollary~\ref{cor:bounds}; \citealp{manamkoong2025}) and extending it to condition (3) of Proposition~\ref{prop:drop}; and a general result on when measuring both groups with the same imprecise instrument beats differential measurement (Section~\ref{sec:nhanes}).

\appendix
\section{Proof of Theorem~\ref{thm:normality}}\label{app:proof}

The doubly robust estimator equals $\theta$ plus an average of i.i.d. mean-zero terms, plus remainders of order $o_P(n^{-1/2})$; the central limit theorem does the rest.

\paragraph{Notation.} We first work with i.i.d. observations $O=(Z,X,Y)$; the case of several observations per unit is treated at the end. $P_nf=n^{-1}\sum_if(O_i)$; $\E f$ is the population expectation; $\|h\|_Q$ and $\|h\|_P$ are the norms of $L^2(Q)$ and $L^2(P)$. Let $\bar o=(1-\eta)/\eta$ and $\bar w=(1-\pi)\bar o/\pi$: by condition 1 and Theorem~\ref{thm:odds}, $o\le\bar o$ and $w\le\bar w$, and by condition 2, $\hat o\le\bar o$. The observations are bounded: $|Y|\le M$ and $|\hat m|\le M$. In fold $k$, with $n_k$ observations, $(\hat o_k,\hat m_k)$ are trained on the data outside the fold; $\E_k$ and $\PP_k$ denote expectation and probability conditional on those data, and $P_{n,k}$ the empirical mean within the fold.

\paragraph{Three identities.} For every bounded measurable $h$:
\begin{align*}
&\text{(I1)}\quad \E[(1-Z)\,o(X)\,h(X)]=\pi\,\E_P[h],\\
&\text{(I2)}\quad \E[(1-Z)\,h(X)\,(Y-m_0(X))]=0,\\
&\text{(I3)}\quad \|h\|_P^2\le\bar w\,\|h\|_Q^2 .
\end{align*}
(I1): $\E[(1-Z)oh]=(1-\pi)\E_Q[oh]$, and $o=\pi w/(1-\pi)$ by Theorem~\ref{thm:odds}, so it equals $\pi\E_Q[wh]=\pi\E_P[h]$. (I2): tower property conditioning on $X$ on $\{Z=0\}$, where $\E[Y\mid X,Z=0]=m_0(X)$. (I3): $\|h\|_P^2=\E_Q[wh^2]\le\bar w\,\E_Q[h^2]$.

\paragraph{Step 1 (the function $\psi$ and its bias).} For bounded $(o',m')$, let $\psi(O;o',m')=(1-Z)\,o'(X)\,(Y-m'(X))+Z\,(m'(X)-\theta)$. Then $\pi\phi=\psi(\cdot;o,m_0)$ and
\[
\E[\psi(O;o',m')]=(1-\pi)\,\E_Q[(o'-o)(m_0-m')].
\]
By (I2) with $h=o'$, $\E[(1-Z)o'(Y-m')]=(1-\pi)\E_Q[o'(m_0-m')]$. By (A2), $\theta=\E_P[m_0]$, and by (I1) with $h=m'-m_0$, $\E[Z(m'-\theta)]=\pi\E_P[m'-m_0]=(1-\pi)\E_Q[o(m'-m_0)]$. Adding gives the formula; in particular, $\E\phi=0$. This is Proposition~\ref{prop:dr} written in odds.

\paragraph{Step 2 (linearization).} With $\hat A=P_n[(1-Z)\hat o(Y-\hat m)]$, $\hat B=P_n[(1-Z)\hat o]$ and $\hat\pi=P_nZ$, the estimator is $\hat\theta=\hat A/\hat B+P_n[Z\hat m]/\hat\pi$. On the event $\{\hat\pi>\pi/2,\ \hat B>\pi/2\}$, whose probability tends to 1 and on which we work from now on, $P_n[Z\theta]/\hat\pi=\theta$ and algebra gives exactly
\[
\hat\theta-\theta=\frac1\pi P_n\psi(O;\hat o,\hat m)+\hat A\Bigl(\frac1{\hat B}-\frac1\pi\Bigr)+P_n[Z(\hat m-\theta)]\Bigl(\frac1{\hat\pi}-\frac1\pi\Bigr).
\]
Steps 3 and 4 show that the first term is $P_n\phi+o_P(n^{-1/2})$; Step 5, that the other two are $o_P(n^{-1/2})$.

\paragraph{Step 3 (empirical process).} Let $\Delta_k=\psi(\cdot;\hat o_k,\hat m_k)-\psi(\cdot;o,m_0)$. Conditional on the data outside fold $k$, $\hat o_k$ and $\hat m_k$ are fixed functions and the $\Delta_k(O_i)$ in the fold are i.i.d. By the conditional Chebyshev inequality, for every $\epsilon>0$,
\[
\PP_k\Bigl(\bigl|(P_{n,k}-\E_k)\Delta_k\bigr|>\epsilon\,n^{-1/2}\Bigr)\le\frac{n}{n_k}\,\frac{\E_k[\Delta_k^2]}{\epsilon^2}.
\]
Since $\Delta=(1-Z)[(\hat o-o)(Y-m_0)+\hat o(m_0-\hat m)]+Z(\hat m-m_0)$ and $|Y-m_0|\le2M$, the inequality $(a+b+c)^2\le3(a^2+b^2+c^2)$ and (I3) give
\[
\E_k[\Delta_k^2]\le3\Bigl[4M^2(1-\pi)\,\|\hat o_k-o\|_Q^2+\bigl((1-\pi)\bar o^{\,2}+\pi\bar w\bigr)\|\hat m_k-m_0\|_Q^2\Bigr],
\]
which goes to 0 in probability by condition 3. Since $n/n_k\to K$, the conditional probability goes to 0 in probability, and being bounded by 1, dominated convergence gives that the unconditional probability also goes to 0. Hence $(P_{n,k}-\E_k)\Delta_k=o_P(n^{-1/2})$ for each of the $K$ folds.

\paragraph{Step 4 (bias).} Since $\E_k\psi(\cdot;o,m_0)=0$, Step 1 and Cauchy--Schwarz give
\[
|\E_k\Delta_k|=(1-\pi)\bigl|\E_Q[(\hat o_k-o)(m_0-\hat m_k)]\bigr|\le(1-\pi)\,r_o\,r_m=o_P(n^{-1/2}).
\]
Summing over folds with weights $n_k/n$, $P_n\psi(\cdot;\hat o,\hat m)=P_n\psi(\cdot;o,m_0)+o_P(n^{-1/2})=\pi P_n\phi+o_P(n^{-1/2})$.

\paragraph{Step 5 (normalizations).} We split the factors of Step 2 into an i.i.d. part and an auxiliary-error part:
\begin{itemize}
  \item $\hat A=P_n[(1-Z)o(Y-m_0)]+P_n[(1-Z)\{(\hat o-o)(Y-m_0)+\hat o(m_0-\hat m)\}]$. The first term is $O_P(n^{-1/2})$ by the central limit theorem. In the second, the centered part is $o_P(n^{-1/2})$ by the argument of Step 3, and its conditional mean, by (I2), is $(1-\pi)\E_Q[\hat o_k(m_0-\hat m_k)]$, bounded by $(1-\pi)\bar o\,r_m$. Hence $\hat A=O_P(n^{-1/2})+O_P(r_m)$.
  \item Likewise, $\hat B-\pi=(P_n-\E)[(1-Z)o]+P_n[(1-Z)(\hat o-o)]=O_P(n^{-1/2})+O_P(r_o)$, using $\E[(1-Z)o]=\pi$ by (I1) and that the $L^1(Q)$ norm does not exceed the $L^2(Q)$ norm.
  \item $P_n[Z(\hat m-\theta)]=P_n[Z(m_0-\theta)]+P_n[Z(\hat m-m_0)]=O_P(n^{-1/2})+O_P(r_m)$, using $\E[Z(m_0-\theta)]=0$ and (I3). Moreover, $1/\hat\pi-1/\pi=O_P(n^{-1/2})$.
\end{itemize}
Since $\hat B\to\pi>0$ in probability,
\[
\hat A\Bigl(\frac1{\hat B}-\frac1\pi\Bigr)=\frac{\hat A(\pi-\hat B)}{\pi\hat B}=O_P\bigl(n^{-1}+n^{-1/2}(r_o+r_m)+r_or_m\bigr)=o_P(n^{-1/2}),
\]
and the third term of Step 2 is $O_P(n^{-1/2})\cdot[O_P(n^{-1/2})+O_P(r_m)]=o_P(n^{-1/2})$.

\paragraph{Step 6 (conclusion).} $\hat\theta-\theta=P_n\phi+o_P(n^{-1/2})$. The $\phi(O_i)$ are i.i.d., with $\E\phi=0$ and $\E\phi^2=V<\infty$, because $\phi$ is bounded. The central limit theorem and Slutsky's lemma give $\sqrt n(\hat\theta-\theta)\to N(0,V)$.

\paragraph{Step 7 (estimating $V$).} The implemented estimator uses
\[
\hat\phi(O)=\frac{(1-Z)\,\hat o(X)\,\bigl(Y-\hat m(X)-\hat A/\hat B\bigr)}{\hat B}+\frac{Z\,(\hat m(X)-\hat C)}{\hat\pi},\qquad\hat C=\frac{P_n[Z\hat m]}{\hat\pi}.
\]
All functions are bounded, so $|\hat\phi-\phi|\le c\,(|\hat o-o|+|\hat m-m_0|+|\hat C-\theta|+|\hat A/\hat B|+|\hat B-\pi|+|\hat\pi-\pi|)$ for a constant $c$ (random, but bounded on that event). The conditional means of the squares of the first two terms go to 0 fold by fold, and the scalars go to 0 in probability, so $P_n(\hat\phi-\phi)^2=o_P(1)$. With $P_n\phi^2\to V$ by the law of large numbers and Cauchy--Schwarz, $|P_n\hat\phi^2-P_n\phi^2|\le P_n(\hat\phi-\phi)^2+2[P_n(\hat\phi-\phi)^2\cdot P_n\phi^2]^{1/2}\to0$. Thus $\hat V\to V$ in probability, and by Slutsky the interval $\hat\theta\pm z_{1-\alpha/2}(\hat V/n)^{1/2}$ has asymptotic coverage $1-\alpha$.

\paragraph{Several observations per unit.} If each unit $i$ contributes a block of $m_i\le\bar m$ observations with the same $Z_i$, units are i.i.d. and $P$ and $Q$ are defined as observation-level laws, all of the above is repeated with block quantities. Identities (I1)--(I3) hold for block sums by linearity, with $\pi$ and $1-\pi$ replaced by $\bar N_1=\E[\sum_jZ_{ij}]$ and $\bar N_0=\E[\sum_j(1-Z_{ij})]$, the expected numbers of target and donor observations per unit (which can exceed 1), and $o$ by the observation-level odds, so that $o=\bar N_1w/\bar N_0$. Cross-fitting by unit keeps the evaluated fold independent of the training data, and $m_i\le\bar m$ keeps everything bounded. The unit-level influence function is $\sum_j\psi(O_{ij};o,m_0)/\bar N_1$, and $\sqrt n$ refers to the number of units; the clustered standard error is its sample version. With dependence within units, $V$ is the variance of this block influence function, not the closed form of Section~\ref{sec:inference}, which is the i.i.d. case.

\begin{thebibliography}{99}\small

\bibitem[Baitairian et~al.(2025)]{baitairian2025} Baitairian, J.-B., Sebastien, B., Jreich, R., Katsahian, S. and Guilloux, A. (2025). Calibrating confounding strength in sensitivity models for weighting estimators: a comparative review and a new method. arXiv:2510.16560.

\bibitem[Bickel et~al.(2009)]{bickel2009} Bickel, S., Brückner, M. and Scheffer, T. (2009). Discriminative learning under covariate shift. \emph{Journal of Machine Learning Research}, 10, 2137--2155.

\bibitem[Blake et~al.(2020)]{blake2020} Blake, H. A., Leyrat, C., Mansfield, K. E., Seaman, S., Tomlinson, L. A., Carpenter, J. and Williamson, E. J. (2020). Propensity scores using missingness pattern information: a practical guide. \emph{Statistics in Medicine}, 39, 1641--1657.

\bibitem[Chernozhukov et~al.(2018)]{chernozhukov2018} Chernozhukov, V., Chetverikov, D., Demirer, M., Duflo, E., Hansen, C., Newey, W. and Robins, J. (2018). Double/debiased machine learning for treatment and structural parameters. \emph{The Econometrics Journal}, 21(1), C1--C68.

\bibitem[Connor Gorber et~al.(2007)]{connorgorber2007} Connor Gorber, S., Tremblay, M., Moher, D. and Gorber, B. (2007). A comparison of direct vs. self-report measures for assessing height, weight and body mass index: a systematic review. \emph{Obesity Reviews}, 8(4), 307--326.

\bibitem[Cook et~al.(2008)]{cook2008} Cook, T. D., Shadish, W. R. and Wong, V. C. (2008). Three conditions under which experiments and observational studies produce comparable causal estimates. \emph{Journal of Policy Analysis and Management}, 27(4), 724--750.

\bibitem[Crump et~al.(2009)]{crump2009} Crump, R. K., Hotz, V. J., Imbens, G. W. and Mitnik, O. A. (2009). Dealing with limited overlap in estimation of average treatment effects. \emph{Biometrika}, 96(1), 187--199.

\bibitem[D'Agostino and Rubin(2000)]{dagostino2000} D'Agostino, R. B., Jr. and Rubin, D. B. (2000). Estimating and using propensity scores with partially missing data. \emph{Journal of the American Statistical Association}, 95, 749--759.

\bibitem[Dahabreh et~al.(2019)]{dahabreh2019} Dahabreh, I. J., Robertson, S. E., Tchetgen Tchetgen, E. J., Stuart, E. A. and Hernán, M. A. (2019). Generalizing causal inferences from individuals in randomized trials to all trial-eligible individuals. \emph{Biometrics}, 75(2), 685--694.

\bibitem[D'Amour et~al.(2021)]{damour2021} D'Amour, A., Ding, P., Feller, A., Lei, L. and Sekhon, J. (2021). Overlap in observational studies with high-dimensional covariates. \emph{Journal of Econometrics}, 221, 644--654.

\bibitem[de~Luna et~al.(2011)]{deluna2011} de Luna, X., Waernbaum, I. and Richardson, T. S. (2011). Covariate selection for the nonparametric estimation of an average treatment effect. \emph{Biometrika}, 98(4), 861--875.

\bibitem[Dehejia and Wahba(1999)]{dehejia1999} Dehejia, R. H. and Wahba, S. (1999). Causal effects in nonexperimental studies: reevaluating the evaluation of training programs. \emph{Journal of the American Statistical Association}, 94(448), 1053--1062.

\bibitem[Dong et~al.(2020)]{dong2020} Dong, J., Zhang, J. L., Zeng, S. and Li, F. (2020). Subgroup balancing propensity score. \emph{Statistical Methods in Medical Research}, 29(3), 659--676.

\bibitem[Dorn and Guo(2023)]{dorn2023} Dorn, J. and Guo, K. (2023). Sharp sensitivity analysis for inverse propensity weighting via quantile balancing. \emph{Journal of the American Statistical Association}, 118(544), 2645--2657.

\bibitem[Hahn(1998)]{hahn1998} Hahn, J. (1998). On the role of the propensity score in efficient semiparametric estimation of average treatment effects. \emph{Econometrica}, 66(2), 315--332.

\bibitem[Heckman et~al.(1997)]{heckman1997} Heckman, J. J., Ichimura, H. and Todd, P. E. (1997). Matching as an econometric evaluation estimator: evidence from evaluating a job training programme. \emph{The Review of Economic Studies}, 64(4), 605--654.

\bibitem[Hirano et~al.(2003)]{hirano2003} Hirano, K., Imbens, G. W. and Ridder, G. (2003). Efficient estimation of average treatment effects using the estimated propensity score. \emph{Econometrica}, 71(4), 1161--1189.

\bibitem[Imai and Yamamoto(2010)]{imai2010} Imai, K. and Yamamoto, T. (2010). Causal inference with differential measurement error: nonparametric identification and sensitivity analysis. \emph{American Journal of Political Science}, 54(2), 543--560.

\bibitem[Khan and Tamer(2010)]{khan2010} Khan, S. and Tamer, E. (2010). Irregular identification, support conditions, and inverse weight estimation. \emph{Econometrica}, 78(6), 2021--2042.

\bibitem[LaLonde(1986)]{lalonde1986} LaLonde, R. J. (1986). Evaluating the econometric evaluations of training programs with experimental data. \emph{The American Economic Review}, 76(4), 604--620.

\bibitem[Liu et~al.(2023)]{liu2023} Liu, Y., Li, H., Zhou, Y. and Matsouaka, R. A. (2023). Average treatment effect on the treated, under lack of positivity. arXiv:2309.01334.

\bibitem[Ma and Wang(2020)]{ma2020} Ma, X. and Wang, J. (2020). Robust inference using inverse probability weighting. \emph{Journal of the American Statistical Association}, 115, 1851--1860.

\bibitem[Ma and Namkoong(2025)]{manamkoong2025} Ma, Y. and Namkoong, H. (2025). A sensitivity approach to causal inference under limited overlap. arXiv:2511.22003.

\bibitem[Ma et~al.(2023)]{ma2023} Ma, Y., Sant'Anna, P. H. C., Sasaki, Y. and Ura, T. (2023). Doubly robust estimators with weak overlap. arXiv:2304.08974.

\bibitem[McClean et~al.(2024)]{mcclean2024} McClean, A., Branson, Z. and Kennedy, E. H. (2024). Calibrated sensitivity models. arXiv:2405.08738.

\bibitem[National Center for Health Statistics(2020)]{nchs2020} National Center for Health Statistics (2020). \emph{National Health and Nutrition Examination Survey, 2017--2018 data files}. Hyattsville, MD: Centers for Disease Control and Prevention. Files DEMO\_J, BMX\_J, WHQ\_J and BPX\_J, \texttt{https://wwwn.cdc.gov/nchs/nhanes}.

\bibitem[Petersen et~al.(2012)]{petersen2012} Petersen, M. L., Porter, K. E., Gruber, S., Wang, Y. and van der Laan, M. J. (2012). Diagnosing and responding to violations in the positivity assumption. \emph{Statistical Methods in Medical Research}, 21(1), 31--54.

\bibitem[Qin(1998)]{qin1998} Qin, J. (1998). Inferences for case-control and semiparametric two-sample density ratio models. \emph{Biometrika}, 85(3), 619--630.

\bibitem[Rassen et~al.(2012)]{rassen2012} Rassen, J. A., Glynn, R. J., Rothman, K., Setoguchi, S. and Schneeweiss, S. (2012). Applying propensity scores estimated in a full cohort to adjust for confounding in subgroup analyses. \emph{Pharmacoepidemiology and Drug Safety}, 21(7), 697--709.

\bibitem[Rudin(1987)]{rudin1987} Rudin, W. (1987). \emph{Real and Complex Analysis} (3rd ed.). McGraw-Hill.

\bibitem[Smith and Todd(2005)]{smith2005} Smith, J. A. and Todd, P. E. (2005). Does matching overcome LaLonde's critique of nonexperimental estimators? \emph{Journal of Econometrics}, 125(1--2), 305--353.

\bibitem[Steiner et~al.(2011)]{steiner2011} Steiner, P. M., Cook, T. D. and Shadish, W. R. (2011). On the importance of reliable covariate measurement in selection bias adjustments using propensity scores. \emph{Journal of Educational and Behavioral Statistics}, 36(2), 213--236.

\bibitem[Sugiyama et~al.(2012)]{sugiyama2012} Sugiyama, M., Suzuki, T. and Kanamori, T. (2012). \emph{Density Ratio Estimation in Machine Learning}. Cambridge University Press.

\bibitem[Tan(2006)]{tan2006} Tan, Z. (2006). A distributional approach for causal inference using propensity scores. \emph{Journal of the American Statistical Association}, 101(476), 1619--1634.

\bibitem[VanderWeele and Shpitser(2011)]{vanderweele2011} VanderWeele, T. J. and Shpitser, I. (2011). A new criterion for confounder selection. \emph{Biometrics}, 67(4), 1406--1413.

\bibitem[Westreich et~al.(2011)]{westreich2011} Westreich, D., Cole, S. R., Funk, M. J., Brookhart, M. A. and Stürmer, T. (2011). The role of the c-statistic in variable selection for propensity score models. \emph{Pharmacoepidemiology and Drug Safety}, 20(3), 317--320.

\bibitem[Westreich et~al.(2017)]{westreich2017} Westreich, D., Edwards, J. K., Lesko, C. R., Stuart, E. and Cole, S. R. (2017). Transportability of trial results using inverse odds of sampling weights. \emph{American Journal of Epidemiology}, 186(8), 1010--1014.

\bibitem[Zhao et~al.(2019)]{zhao2019} Zhao, Q., Small, D. S. and Bhattacharya, B. B. (2019). Sensitivity analysis for inverse probability weighting estimators via the percentile bootstrap. \emph{Journal of the Royal Statistical Society: Series B}, 81(4), 735--761.

\end{thebibliography}

\end{document}
